% Stage I: the C program's choice of the sieving moduli and the cut points of its stages by the
% curve, the blocks of the run with the correction of the moduli between them, the numbers of the
% moduli used (message 4), the report of the library under -v, the tuning options -r -R -U -C -A -Q
% -W of the command-line program, and the block and lag lemmas. Proposed names are set with \lname
% and become \lean once the declarations exist.
\chapter{The choice of the moduli and the blocks}\label{chap:moduli}

Stage~I models what is left of the C code of ratpoints 3.1.1 apart from its threads: how the library
chooses the moduli of its sieve and the cut points of its three stages from the curve, and how it
corrects that choice while it runs. \verb|sieving_info| examines the primes of the table and rates
each by its density (the share of the numerators that its rows let through); ranks them by keys
(what a modulus costs per word swept, divided by $-\log r$ for its density $r$); adds composite
moduli whose tables the run can pay for; ends the first stage by a rule that weighs the survivors
that the next modulus removes against its cost; adds, for the second stage, an offset that grows
with the length of the run; chooses the third stage by a rule of its own, which may look at further
primes; and keeps the primes left as a reserve. All of this is arithmetic in doubles on an estimate
of the run, \verb|run_shape|, from the shape of the loop over the denominators, the search region
and the packing of the classes. The run then sifts its denominators in blocks of consecutive steps
of that loop; after each block it folds counters of the sieve into a state from which
\verb|adapt_primes| may move the ends of the second and the third stage, a decision that takes
effect $\mathrm{RP\_ADAPT\_LAG} + 1 = 33$ blocks after the block folded. The search writes back the
numbers of the moduli that its last block used, which the program prints after the points
(message~4), and the library prints a report of its own under \verb|-v|, from the arguments as it
normalizes them to the moduli chosen and the forbidden divisors. The command-line program reads the
tuning options \verb|-r|, \verb|-R|, \verb|-U|, \verb|-C|, \verb|-A|, \verb|-Q| and \verb|-W| into
the arguments of the library.

Stage~I models all of this as the C code has it, on the sequential path, to which the pool of
threads is equivalent (Theorem~\ref{thm:threads}): the doubles of the cost model are values of Lean
core's \verb|Float|, IEEE~754 binary64 with the arithmetic of the platform, and every expression
keeps the order of the operations of the C code; the sorts are stable merge sorts with the
comparisons of the C code; the blocks, the counters, the decisions, the folds and the correction are
those of the C code, with the counters as $64$-bit words that wrap as the C code's
\verb|unsigned long|s do. The listing of the denominators of a shape, which the model had in place
of the blocks since stage~E, goes from the model; the proofs keep it as a definition of their own.
The choice and the run produce exactly the moduli, the cut points, the decisions and the counters of
the C program on the machine on which both run. After stage~I every line of code of the C files at
the pin is modelled by a declaration of the model or left out with a reason recorded in the
correspondence, stage~I adding the threads, a probe of the development, a macro that the C code
reads nowhere and the buffers of memory of the arguments.

For every input the points that $\mathrm{findPoints}$ returns do not change: a modulus removes only
candidates that are not points, whatever moduli the choice takes and whatever decisions the run
makes, and the blocks visit the denominators of the shape in the order of its loop. What changes is
what depends on the choice: the survivors under \verb|RATPOINTS_NO_CHECK| (\verb|-x|), message~4 and
the report under \verb|-v|, which are now those of ratpoints, byte for byte. The interface of the
search changes with them. $\mathrm{findPoints}$ and $\mathrm{findPointsWith}$ take the options of
the sieve as the caller of the library gives them, the fields of \verb|ratpoints_args| that tune the
sieve (Definition~\ref{def:sieveOptions}), now with the values of the tuning options and the array
size, in place of explicit parameters and their well-formedness; the result of the choice,
$\mathrm{Params}$ (Definition~\ref{def:Params}), is internal to the search. So the statements that
comparator certifies drop the parameters $P$ and the hypothesis $P.\mathrm{WF}$: the six on the
points and the one on the protocol of the delivery take the options $s$ of the sieve with no
hypothesis on them, the two refusal theorems take $s$ in place of $P$, and the theorem that the
parameters of the program are well-formed goes, since nothing is left for it to say. Eleven
statements remain: the two on the brute-force reference unchanged, and the nine on the search each
stronger than before, since they cover the program's own choice of the moduli, whatever the values
of the doubles of its cost model, and every option of the sieve. The specification that the
challenge restates does not change.

The theorems say nothing about the values of the doubles, and need not (Section~\ref{sec:costs}):
every contract of the choice, of the blocks and of the correction holds for any values that the cost
model computes, so that the cut points and the decisions are, for the proofs, arbitrary within the
caps that the natural numbers of the C code impose. What the proofs show is that the choice is what
the run needs (every modulus sound, the tables and the store complete, the entries from the end of
the second stage on primes of the table, the forbidden divisors sound), that the correction keeps
every decision so, and that the run of blocks delivers, in the order of the loop of the shape, the
points of each denominator with the moduli of the decision of its block. Since the counters are
counted before the exact check (Section~\ref{sec:moduliC}), the decisions depend neither on the
callback nor on the flags of the output; the proofs give the state of the correction in closed form,
as a fold of the reports of the blocks, and with it a sound test of the numerators that the
certified statements need, the same for every callback. Three further theorems are stated and proved
in this chapter but not certified: the block lemma, that the run delivers the points of its blocks
in order, each a function of its block and the decision of the block (Theorem~\ref{thm:blocks}); the
lag lemma, that the decision of a block reads only the folds of the blocks at least
$\mathrm{RP\_ADAPT\_LAG} + 1$ before it (Theorem~\ref{thm:lag}); and, from both, that blocks sifted
apart and delivered in order give what the run delivers (Theorem~\ref{thm:threads}), the property by
which the threaded path of the C code, which sieves blocks apart and delivers them in order, gives
the output of the sequential path.

The doubles bring one dependence on the machine: glibc's \verb|log| and \verb|pow| are not correctly
rounded, and glibc chooses their variant by the processor (Section~\ref{sec:costs}), so that the
choice of ratpoints depends on the machine it runs on, and the model, which calls the same functions
of the platform's C library, is exact relative to that library; the differential tests compare the
two programs on one machine. As in the other stages, since the aim of the project is to verify the C
code as far as that can be done, the model follows the data, the loops and the tests of the C code
where the C code has them, also in two points of the choice that look unintended: the sort of the
third stage by density starts at the end of the second stage, not of the third, so that the third
stage can sieve with a prime that its rule did not take, and \verb|RATPOINTS_DEFAULT_SP1|
(\csource{ratpoints.h}{42}{42}) is read nowhere. Both are recorded with the C code as cosmetic notes
for a release after 3.1.1; the C code stays at 3.1.1 while the project runs, so that the project
certifies that release, and changes to it come later with the adaptations here. The speed of the
model is not an aim.

\section{What the C code does}\label{sec:moduliC}

\paragraph{The tuning of the choice.} The arguments of the library hold, besides the number of
primes to look at and the cut points, the tuning fields of the choice
(\csource{ratpoints.h}{365}{367}): the target \verb|survivors_per_word| of the first stage, the
offset \verb|sp2_extra| of the second stage in a long run and its scale \verb|sp2_u0|, the cost
\verb|cost_table| of a row of a table, the mode \verb|adapt| of the correction, the number
\verb|sp3_extra| of primes of the third stage, which the model has read since stage~F2, the cost
\verb|sp3_per_denom| of a prime of the third stage per denominator, and the cost \verb|check_cost|
of one exact check. \verb|main.c| sets them to $-1$ and $-1.0$, the defaults
(\csource{main.c}{337}{347}), and reads \verb|-r -U -C -Q -W| as doubles with \verb|scan_double| and
\verb|-R -A| as \verb|long|s, as it has read \verb|-P| since stage~F2. \verb|sieving_info| takes the
target of the first stage when it is positive and the offset and the scale of the second stage and
the cost of a row when they are at least $0$, and the default macro of \verb|ratpoints.h| otherwise
(\csource{find_points.c}{2569}{2575}), so that a NaN of a library caller gives the default; the
other fields are read where they are used: the cost of a check when it is positive, and otherwise
computed from the curve (line 1293); the cost of a prime of the third stage per denominator when it
is at least $0$, and otherwise its macro (lines 1489 and 2929); the number of primes of the third
stage when it is at least $0$ and below the primes left, a negative one meaning the cost rule of the
stage (lines 2926-2928, 2974 and 1486); and \verb|adapt| is $0$ for no correction, $1$ for the third
stage only, $2$ or more for the second stage too, and negative for $1$ (lines 1395 and 2580). The
number of primes to look at is the normalized \verb|num_primes|, and whether the caller left it to
the library, \verb|np_is_default|, is read before the normalization
(\csource{find_points.c}{4577}{4577}).

\paragraph{The densities.} \verb|examine_prime| also counts the residues $a$ for which $f(a)$ is a
square modulo $p$, zero included, $n_p$, and returns the code $-1$ for no such residue and no point
at infinity modulo $p$ (the exit), $0$ when every residue passes and there is no point at infinity
(the prime says nothing), $2$ when every residue passes otherwise (the prime says only that
numerator and denominator are coprime) and $1$ otherwise, with the density
$r = (n_p + 1)(p - 1)/p^2$ when there is a point at infinity modulo $p$ and $r = n_p / p$ otherwise
(\csource{find_points.c}{1766}{1781}). \verb|examine_power| computes the density of a prime power
$p^e$ from its tables, with the contributions of the multiples of $p$ when denominators divisible by
$p$ occur, and calls the power useful when its density is below $1 - 10^{-9}$
(\csource{find_points.c}{2306}{2331}). The sieve entries keep the density and the count of their
factors, the forbidden divisors the reciprocal of their prime.

\paragraph{The cost model.} The candidates of the ranking are records with their density, key, cost,
modulus, the mask of the indices of their primes, their sieve entry and the codes of their factors
(\csource{find_points.c}{147}{149}); two comparisons order them by key and by density
(\csource{find_points.c}{1194}{1206}). The cost of a modulus per word swept, in the units of an
\verb|and| of the first stage (\verb|prime_cost|), its key (\verb|prime_key|), the key of the first
stage (\verb|phase_1_key|), the cost of one exact check from the bit lengths of the coefficients and
the height bound (\verb|check_cost|), the factor of what a surviving word costs downstream
(\verb|downstream_factor|), the rule that ends the first stage (\verb|phase_1_wants|,
\verb|primes_for_phase_1|) and the offset of the second stage (\verb|phase_2_offset|) are
\csource{find_points.c}{1220}{1318} and \csource{find_points.c}{1539}{1626}. \verb|run_shape|
(\csource{find_points.c}{1836}{2106}) estimates the run before any sieving: for each segment of
steps of the loop over the denominators (see the blocks below) it samples $64$ points of the segment
of $k$, or $64$ denominators, counts the numerators of the search intervals that each sampled
denominator considers (\verb|numerators_for|), the classes kept and what the rounding of an interval
cut by the height bound to whole bit arrays adds (\verb|clipped_padding|), the share of the
denominators that the forbidden divisors leave (\verb|forbidden_fraction| for the valuations) and
the Jacobi factor $0.53$, and returns the denominators $D$, the words $U$, the padding among them
and the calls of the sieve, $D$, $U$ and the calls floored at $1$ and the padding $0$ when $U$ is.
\verb|mean_bits_per_word| (\csource{find_points.c}{2128}{2159}) weighs the bits per word of the
first word of each class kept by its share of the numerators.

\paragraph{The loop of the tables.} \verb|sieving_info| (\csource{find_points.c}{2522}{3103}) first
estimates the run without forbidden divisors and derives the cost of the calls per word, the bits
per word swept and the offset of the second stage; then it examines the primes of the table in turn
(\csource{find_points.c}{2601}{2707}): at the code $-1$ it returns the prime, and the search returns
with no point; a prime that says something becomes a candidate with its key for the first stage; the
forbidden divisor of the prime is made as before. When the caller left the number of primes to the
library and the last prime to look at has been looked at, it sorts the candidates by key, computes
how many the first stage would take and the second after it, and looks at one more prime while the
candidates that say more than coprimality are fewer (\csource{find_points.c}{2693}{2705}), at most
$\mathrm{numPrimes}$. After the further forbidden divisors and the switch of the checked loop it
estimates the run again, with the forbidden divisors, and keys the candidates again.

\paragraph{The composite moduli.} \verb|add_moduli| (\csource{find_points.c}{2346}{2411}) offers the
odd composite $m$ from $9$ up to \verb|RATPOINTS_COMPOSITE_MAX| and below
$2^{\mathrm{maxBitsInPrime}}$ whose table the run can pay for, $\mathrm{ct}\, m \min(D, m)/U \le 2$
for the cost $\mathrm{ct}$ of a row of a table, that factor over the primes looked at into at most
\verb|RP_MAX_FACTORS| factors among the first $64$ primes, each prime with an entry and each prime
power useful, its density the product of those of its factors; the prime powers are examined once
each. \verb|make_modulus| (\csource{find_points.c}{2414}{2475}) makes the sieve entry of a composite
modulus that the ranking takes, from the codes of its factors.

\paragraph{The first two stages.} The candidates are sorted by key; the first stage takes them in
that order (\verb|take_entries|, \csource{find_points.c}{2490}{2520}), passing over and deleting a
candidate that shares a prime with one taken, $s_1$ of them when \verb|-n| gives $s_1$ and otherwise
while the rule of the first stage wants the next; the rest is keyed again for the second stage, by
the cost of a word that survives the first stage, and sorted again; the second stage takes
$s_2 - s_1$ more when \verb|-N| gives $s_2$ and otherwise the offset
(\csource{find_points.c}{2843}{2871}). The candidates after the second stage that are not primes
saying more than coprimality are deleted (\csource{find_points.c}{2879}{2890}), and the entries of
the composite moduli taken are made (\csource{find_points.c}{2894}{2900}).

\paragraph{The third stage and the reserve.} The third stage (\csource{find_points.c}{2922}{3016})
takes the prime of least density left, the first on ties, up to as many primes beyond $s_2$ as
\verb|-P| gives, and without \verb|-P| up to $\mathrm{numPrimes}$ while the survivors per
denominator that it removes are worth its cost per survivor and per denominator. When the caller
left the number of primes to the library, it may look at one more prime when no candidate is left
or, without \verb|-P|, when none is worth it and a prime of the mean density of the candidates that
say more than coprimality would still be worth it (lines 2973-2992); that prime can end the search
when the curve has no points modulo it (\csource{find_points.c}{2994}{2998}). Then it sorts the
candidates from the position $s_2$ on by density (\csource{find_points.c}{3024}{3030}), the third
stage with the primes left, which are the reserve of the correction, and keeps the sieve entries and
the reciprocals of all of them (\csource{find_points.c}{3064}{3070}); the sort from $s_2$, not from
the end $s_3$ of the third stage, is the first of the two cosmetic points of the introduction. It
records the end $s_3$, the number of the entries, the estimates of the run, the cost of a check
relative to \verb|RATPOINTS_CHECK_REFERENCE| and the number of words at which the first correction
is made, $10^6$ when the mode of the correction is not $0$ and neither \verb|-n| nor \verb|-N| is
given, and never otherwise (\csource{find_points.c}{2580}{2588}).

\paragraph{The blocks.} \verb|run_setup| cuts the loop of the shape into segments of steps
(\csource{find_points.c}{3812}{3837}): the $k$ of the squares $k^2$, for the divisor shape one
segment for each divisor, the words of $64$ denominators, or the denominators themselves. The blocks
have $\lfloor \mathrm{steps}/\mathrm{RP\_BLOCKS\_WANTED} \rfloor + 1$ steps, at least
\verb|RP_BLOCK_MIN| $= 16$, or the length \verb|RP_BLOCK_LEN| when the build defines it, and each
segment holds $\lfloor (\mathrm{hi} - \mathrm{lo})/\mathrm{len} \rfloor + 1$ of them
(\csource{find_points.c}{3838}{3850}). The first decision is that of the choice, for block $0$
(\csource{find_points.c}{3853}{3863}). \verb|run_block| (\csource{find_points.c}{3889}{4078}) finds
the segment and the steps of block $i$, the last decision whose first block is at most $i$, resets
the counters of the sieving thread, and runs the loop of the shape over the steps of the block, with
the ends $s_2$ and $s_3$ of that decision; in the shape of the words the positions of the forbidden
divisors are set at the first word of each block (\csource{find_points.c}{3975}{3988}).
\verb|worker_stat| copies the counters with the decision into the report of the block
(\csource{find_points.c}{4082}{4093}).

\paragraph{The counters.} The sieving thread counts six things, as \verb|unsigned long|s
(\csource{rp-private.h}{579}{581}): the words swept by the first stage (\csource{sift.c}{548}{548}),
the bit arrays that survive it (line 1092 of \verb|sift.c|), the numerators extracted (lines 1166,
1233 and 1271), those that pass the test for common factors and those that pass the third stage
(lines 441 and 444), and the calls of \verb|sift| (\csource{find_points.c}{3461}{3461}). All of them
are counted before the exact check, so that they depend neither on the callback nor on the flags of
the output.

\paragraph{The correction.} \verb|run_fold| (\csource{find_points.c}{4102}{4127}) folds the report
of the next block into the counters of the run, restarting those from the extracted numerators on
when the block was sifted with another end $s_2$ than they were counted under; once the words swept
reach the words of the next correction, it calls \verb|adapt_primes|
(\csource{find_points.c}{1383}{1515}), which sets the words of the next correction to twice the
words swept so far, wrapping as an \verb|unsigned long|, and, from the rates that the run has shown,
the bit arrays per word after the first stage and the numerators per word after the second, fits the
survivors per word as a linear function of the product of the densities, may walk the end of the
second stage up through the entries beyond it (the third stage, then the reserve) or down through
the primes below it (mode $2$ and more), and, without \verb|-P| and once denominators have been
sifted, recomputes the end of the third stage by its rule from the new $s_2$, and otherwise raises
it to $s_2$ when it is below; when either end changed, \verb|run_fold| appends a decision for the
blocks from $j + \mathrm{RP\_ADAPT\_LAG} + 1$ on, $j$ the index of the block folded. The walk down
stops at a composite modulus and above $s_1$, so that the entries from every end $s_2$ on are primes
of the table. The sequential loop (\csource{find_points.c}{5152}{5160}) runs, folds and stops after
each block in turn, and the ends $s_2$ and $s_3$ of the last block folded are written back
(\csource{find_points.c}{5163}{5163}).

\paragraph{The threads.} With threads (\verb|-t|) a pool sieves the blocks apart
(\csource{find_points.c}{4129}{4519}): a block starts only when the fold of the block
$\mathrm{RP\_ADAPT\_LAG} + 1$ before it is done, so that every decision for it is published, each
block is sifted by one thread with the decision in force, and the calling thread delivers the points
of the blocks in their order and stops at the first stop of the callback, as the sequential loop
does. The tables of rows are built under a lock and depend only on their residue. So the output is
that of the sequential loop for every number of threads.

\paragraph{What the search writes back, and message 4.} \verb|find_points_work| zeroes the numbers
of the moduli used (\csource{find_points.c}{4551}{4552}); after the run the search writes the end
$s_1$ of the choice and the ends $s_2$ and $s_3$ of the last block folded
(\csource{find_points.c}{5210}{5213}), and every return before the sieve leaves $0$, $0$, $0$.
\verb|main.c| prints them after the points, unless \verb|-q| or \verb|-v|, after two line feeds and
before messages~7 and~8, as the moduli of the first stage, of both stages together and the further
primes of the third stage (\csource{main.c}{992}{998}).

\paragraph{The report under -v.} Under \verb|-v| the library prints a report of its own on the
standard output, between the points that the callback of the program prints there: the arguments as
it normalizes them (\csource{find_points.c}{4691}{4725}), the reversal check, the Sturm step, the
set-up of the denominators, the packing of the classes, the number of primes to look at, the bits
per word and the moduli chosen with the cut points (\csource{find_points.c}{3089}{3095}), the lists
of the moduli of the three stages, the cost of one exact check and the forbidden divisors
(\csource{find_points.c}{5027}{5066}), the points at infinity and the sieve, with \verb|%.0f|,
\verb|%.1f| and \verb|%.2f| for its doubles and \verb|%lf|, six places, for the intervals of the
Sturm step (line 4875). \verb|main.c| sets the flag (line 583), clears it under \verb|-q| (line 688)
and under \verb|-v| leaves out messages~4, 6, 7 and~8.

\paragraph{The doubles of the C code.} \verb|find_points.c| is compiled with \verb|-O2 -mavx2|,
without \verb|-ffast-math| and without fused multiply-adds, so that every operation on doubles is
the correctly rounded IEEE~754 binary64 operation in the order of the source; \verb|sqrt| is the
instruction, correctly rounded, and \verb|ceil| and \verb|floor| are exact instructions; a
\verb|long| converts to a double by \verb|vcvtsi2sd| and an \verb|unsigned long| counter by the
correctly rounded sequence of gcc, both rounding to nearest with ties to even and exact below
$2^{53}$, and the one conversion of a double to a \verb|long|, \verb|(long)(e + 0.5)| in
\verb|phase_2_offset| (line 1625), truncates a value in range; gcc folds \verb|log(2.0)| at compile
time (line 1289) into the value of glibc; \verb|log| and \verb|pow| are glibc's, which are not
correctly rounded (on the machine of the tests $21$ of $62230$ arguments of \verb|log| and $24$ of
$30000$ of \verb|pow| gave a result one unit in the last place away from the correctly rounded one,
among them arguments that occur, such as the density $111/127$ and $H + 1$ for some heights) and of
which glibc chooses a variant by the processor. The keys are finite or $+\infty$, never a NaN, since
\verb|run_shape| floors $D$ and $U$ at $1$; on such keys the comparison of the sort is a total
preorder, and glibc's \verb|qsort|, a stable merge sort when it can allocate its buffer, gives the
order of every stable sort. Ties occur at finite keys, under \verb|-C 0| or a subnormal \verb|-C|,
where two moduli of equal density get equal keys (such as $15$ and $45$, or $7$ and $49$), and at
$+\infty$ under huge values of \verb|-C|; on all of them a stable merge sort gives the order of
glibc's \verb|qsort|, since each sort takes its input in the order of the C code.

\section{The constants}\label{sec:moduliConstants}

\begin{definition}[The constants of the choice and of the blocks]\label{def:choiceConstants}
  \uses{def:barrett, def:primeTable}
  \lean{Ratpoints.survivorsPerWord, Ratpoints.sp2U0, Ratpoints.sp3PerSurvivor, Ratpoints.sp3PerDenom, Ratpoints.sp3Coprime, Ratpoints.costTable, Ratpoints.compositeMax, Ratpoints.costBp, Ratpoints.costSetup, Ratpoints.costCall, Ratpoints.costLine, Ratpoints.costPhase2, Ratpoints.costSurvivor, Ratpoints.costCheck, Ratpoints.checkStep, Ratpoints.checkLimb, Ratpoints.checkCall, Ratpoints.checkRoot, Ratpoints.checkReference, Ratpoints.maxModuli, Ratpoints.maxFactors, Ratpoints.modulusTableMax, Ratpoints.runShapeSamples, Ratpoints.invStride, Ratpoints.adaptWords, Ratpoints.adaptArrays, Ratpoints.adaptBits, Ratpoints.adaptLag, Ratpoints.blocksWanted, Ratpoints.blockMin, Ratpoints.blockLen, Ratpoints.blockLen_pos, Ratpoints.sp2Extra, Ratpoints.defaultNumPrimes, Ratpoints.defaultMaxForbidden}
  \leanok
  The tuning macros of \verb|ratpoints.h| that the choice reads, each a definition of the model
  citing its macro, \verb|Float|-valued where the C code's value is a double and a natural number
  otherwise: the target of the first stage $0.002$ (\csource{ratpoints.h}{92}{94}), the offset of
  the second stage $\mathrm{sp2Extra} = 11$ (\csource{ratpoints.h}{95}{97}) and its scale
  $1.6 \cdot 10^6$ (\csource{ratpoints.h}{120}{122}), the costs of a prime of the third stage per
  survivor and per denominator and the share of the survivors that are coprime
  (\csource{ratpoints.h}{150}{152}, \csource{ratpoints.h}{153}{155},
  \csource{ratpoints.h}{164}{166}), the cost of a row of a table (\csource{ratpoints.h}{192}{194}),
  the largest composite modulus offered, $64$ (\csource{ratpoints.h}{210}{212}), the costs of an
  entry of the list of residues, of the set-up of a modulus, of a call of the sieve, of a line of
  the walk, of a word of the second stage, of a survivor and of a check
  (\csource{ratpoints.h}{213}{218}, \csource{ratpoints.h}{224}{226},
  \csource{ratpoints.h}{241}{243}, \csource{ratpoints.h}{258}{260}, \csource{ratpoints.h}{261}{263},
  \csource{ratpoints.h}{270}{272}, \csource{ratpoints.h}{280}{282}), the parts of the cost of an
  exact check and its reference (\csource{ratpoints.h}{312}{314}, \csource{ratpoints.h}{315}{317},
  \csource{ratpoints.h}{318}{320}, \csource{ratpoints.h}{321}{323},
  \csource{ratpoints.h}{329}{331}), and the defaults of the number of primes to look at, $30$, and
  of the bound of the forbidden divisors, $64$ (\csource{ratpoints.h}{333}{340},
  \csource{ratpoints.h}{345}{352}); and the macros of the C files: the bound on the number of
  composite moduli offered, half of $\mathrm{maxPrimeEven}$ (\csource{find_points.c}{153}{153}), the
  most factors of a composite modulus, $3$ (\csource{rp-private.h}{520}{520}), the bound $2.0$ of
  the cost of the table of a composite modulus (\csource{find_points.c}{2337}{2337}), the $64$
  samples of the estimate of the run and the reciprocals of the strides
  (\csource{find_points.c}{1917}{1917}, \csource{find_points.c}{1921}{1922}), the words, bit arrays
  and numerators that the correction needs, $10^6$, $1000$ and $200$
  (\csource{find_points.c}{1354}{1354}, \csource{find_points.c}{1355}{1355},
  \csource{find_points.c}{1356}{1356}), the lag $32$ of a decision
  (\csource{find_points.c}{1379}{1381}), the number of blocks wanted, $1024$, and the least length
  of a block, $16$ (\csource{find_points.c}{3729}{3731}, \csource{find_points.c}{3732}{3734}). The
  length of the blocks that a build may fix, \verb|RP_BLOCK_LEN|, is a constant of type
  \verb|Option Nat|, none in the default build, with the decided fact that a length it fixes is
  positive, which the loops over the steps of a block need to end. The offset and the two defaults
  are the constants of the parameters of the executable (Definition~\ref{def:sieveParams}), which
  move here with their links at I.1, so that the options of the sieve
  (Definition~\ref{def:sieveOptions}) read them from this node.
\end{definition}

As for the constants of the earlier stages, no proof depends on their values: every proof holds for
any values of the doubles, and of the new natural numbers the proofs use only that a fixed block
length is positive. The build with other values of the constants learns \verb|Float| constants and
the option, and its variants include the lag $0$ and $1$, the block lengths $1$ and $3$ and one
block wanted, never a block length $0$. A \verb|#guard| checks the bits of every \verb|Float|
literal against those that gcc compiles for the macro. \verb|RATPOINTS_DEFAULT_SP1|
(\csource{ratpoints.h}{42}{42}), which the model had as $\mathrm{defaultSp1}$ for the parameters of
the executable, goes with them at I.12: the C code reads it nowhere.

\section{The densities and the entries}\label{sec:densities}

\begin{definition}[The densities]\label{def:densities}
  \uses{def:powerWords, def:examinePrime}
  \lean{Ratpoints.Sieve.powerDensity, Ratpoints.Sieve.setRevBits}
  \leanok
  The density of a prime power $p^e$ from its tables: the density of the residues and, when
  denominators divisible by $p$ occur, of the multiples of $p$ by their valuations, in the order of
  the C code (\csource{find_points.c}{2306}{2329}); and the loop that sets the bits of \verb|gsq|,
  the multiples $t$ of $p$ below $p^e$ at which the reversed form $\tilde f(t)$ is a square modulo
  $p^e$, split off (\csource{find_points.c}{2294}{2300}). With them the components of the earlier
  chapters gain the densities and the fields that the choice and the correction read. The
  examination of a prime (Definition~\ref{def:examinePrimeExit}) returns, beside the table and the
  coefficients modulo $p$, the C code's code ($0$, $1$ or $2$; at the code $-1$, the exit, it
  returns none, as before) and the density, counting $n_p$ in its loop over the residues; its link
  becomes \verb|find_points.c:1761-1781|. The examination of a prime power
  (Definition~\ref{def:examinePower}) takes whether denominators divisible by $p$ occur and returns
  the density with its tables, and the power is useful when its density is below $1 - 10^{-9}$
  (\lname{Ratpoints.Sieve.Power.useful}, line 2330, which joins that node); its links become
  \verb|find_points.c:2267-2282|, \verb|find_points.c:2285-2293|, \verb|find_points.c:2301|,
  \verb|find_points.c:2330| and \verb|find_points.c:2331|. A sieve entry
  (Definition~\ref{def:SieveEntry}) gains its density $r$ and the count \verb|nf| of its factors
  ($0$ for a prime, $1$ for a prime power, the number of factors of a product), which the correction
  reads; a forbidden divisor tested with words (Definition~\ref{def:forbidden}) gains the reciprocal
  of its prime, with which each block starts its words.
\end{definition}

\begin{lemma}[The entries with their densities]\label{lem:densities}
  \uses{def:densities, lem:power_admits, lem:sieveEntry_admits, def:forbidden}
  \lean{Ratpoints.Sieve.admits_run_examinePower, Ratpoints.Sieve.IsForbiddenEntry.magic_eq, Ratpoints.Sieve.isSound_prime_examinePrime, Ratpoints.Sieve.isSound_power_run_examinePower, Ratpoints.isModulusEntry_of_forall_isSound, Ratpoints.inverses_mulMod_rbaLength_of_isSound, Ratpoints.one_lt_and_coprime_rbaLength_mod, Ratpoints.mul_rbaLength_mod_and_lt, Ratpoints.rbaLength_mul_lt_two_pow_64}
  \leanok
  The contracts of the examinations keep what they said, on the parts of the results that they read:
  the table and the coefficients of a prime, with the count $n_p$ of the true entries of the table
  as part of the invariant; the rows of a prime power, which do not depend on whether denominators
  divisible by $p$ occur. A forbidden divisor tested with words has the reciprocal of its prime. The
  soundness of a factor of a modulus holds for every modulus that the factor divides: a prime factor
  for $p \mid m$, a prime power factor for $p^e \mid m$, with $1 \le e$; and an odd modulus
  $1 < m < \mathrm{maxPrimeEven}$ whose factors are sound and multiply to $m$ has a sound sieve
  entry, whatever its density, count, reference and builder, with the inverse of
  $\mathrm{rbaLength}$ modulo a prime read from its table of inverses at the remainder by the
  reciprocal, and the products of the packing within $64$ bits.
\end{lemma}

\begin{proof}
  \uses{def:forbiddenValuations, def:sieves0, lem:admits_examinePrime, lem:admits_mul, lem:fetchRow,
    lem:sieves0, lem:stage3}
  \leanok
  The contracts of the examinations take the new parts as values: the density is a \verb|Float| that
  no proof reads, and the count of the residues is the number of true entries of the table, a
  conjunct of the invariant of the loop. The soundness of the entries of composite moduli is that of
  Lemma~\ref{lem:sieveEntry_admits}, freed from the factorization of the explicit route: it reads
  only the divisibility of the modulus by each factor and their product.
\end{proof}

\section{The cost model}\label{sec:costs}

\begin{definition}[The keys and the rules of the stages]\label{def:costModel}
  \uses{def:choiceConstants, def:SieveEntry, def:rbaPack}
  \lean{Ratpoints.ModulusCandidate, Ratpoints.ModulusCandidate.keyLe, Ratpoints.ModulusCandidate.rLe, Ratpoints.sortFrom, Ratpoints.checkCost, Ratpoints.primeCost, Ratpoints.primeKey, Ratpoints.phase1Key, Ratpoints.downstreamFactor, Ratpoints.phase1Wants, Ratpoints.primesForPhase1, Ratpoints.phase2Offset}
  \leanok
  A candidate modulus of the ranking: its density, key and cost, the modulus, the mask of the
  indices of its primes (a prime from the $65$th of the table on is in no mask), its sieve entry,
  made for a prime by its examination and for a composite modulus when the ranking takes it, an
  option for the null pointer, the count of its factors ($0$ for a prime, $-1$ for a prime that says
  only coprimality, the number of factors of a composite modulus) and the codes of its factors, the
  index of a prime or $\mathrm{numPrimes}$ plus the index of an examined prime power
  (\csource{find_points.c}{147}{149}); the orders by key and by density, as the test of a merge sort
  whether the comparison of the C code is at most $0$ (\csource{find_points.c}{1194}{1198},
  \csource{find_points.c}{1202}{1206}); and the stable sort of the candidates from a position on.
  The cost of one exact check, the value of \verb|-W| when positive and otherwise from the bit
  lengths of the coefficients and of the height bound (\csource{find_points.c}{1285}{1318}); the
  cost of one more modulus per word swept (\csource{find_points.c}{1220}{1235}); its key, the cost
  over $-\log r$, and $10^{300}$ when that is not positive (\csource{find_points.c}{1241}{1247});
  the key of the first stage, with the cost of the walk (\csource{find_points.c}{1256}{1270}); the
  factor of what a word that survives costs downstream, with \verb|pow| for the mean density of the
  next candidates (\csource{find_points.c}{1539}{1561}); the rule that takes the next candidate into
  the first stage (\csource{find_points.c}{1583}{1589}) and the number of candidates that the first
  stage takes by it (\csource{find_points.c}{1598}{1609}); and the offset of the second stage, the
  offset of a long run scaled by $1/(1 + u_0/U)$ and rounded, at most
  $\mathrm{numPrimes} + \mathrm{maxModuli}$ (\csource{find_points.c}{1619}{1626}).
\end{definition}

\begin{definition}[The estimate of the run]\label{def:runShape}
  \uses{def:numClasses, def:Double, def:ctz, def:forbidden, def:arraySize, def:choiceConstants}
  \lean{Ratpoints.numeratorsFor, Ratpoints.clippedPadding, Ratpoints.forbiddenFraction, Ratpoints.ShapeArgs, Ratpoints.ShapeAcc, Ratpoints.RunEstimate, Ratpoints.sampleSteps, Ratpoints.addClass, Ratpoints.shapeSquares, Ratpoints.sampleDivisor, Ratpoints.shapeSquares1, Ratpoints.keepChecked, Ratpoints.shapeAll, Ratpoints.shapeEstimate, Ratpoints.runShape, Ratpoints.popcount, Ratpoints.bpwAdd, Ratpoints.bpwSquares1, Ratpoints.meanBitsPerWord}
  \leanok
  The numerators that a denominator considers, over the search intervals, whose ends enter the
  arithmetic as \verb|Float|s by their bits (\csource{find_points.c}{1836}{1854}); what the rounding
  to whole bit arrays adds at an end cut by the height bound (\csource{find_points.c}{1863}{1876});
  the share of the denominators whose valuation at $p$ a mask of a forbidden divisor excludes
  (\csource{find_points.c}{1926}{1937}); what the estimate reads of the search (the shape, the
  Jacobi switch, the height bound, the range of the denominators, the search intervals, the word and
  the packings of the classes, the array size), its accumulators with their initial values
  (\csource{find_points.c}{1945}{1957}) and its result; the samples of a segment, written once for
  the squares and for every denominator (\csource{find_points.c}{1964}{1973},
  \csource{find_points.c}{2028}{2037}); a class kept (\csource{find_points.c}{1980}{1981},
  \csource{find_points.c}{2017}{2018}, \csource{find_points.c}{2045}{2046}); the accumulators of the
  squares (\csource{find_points.c}{1960}{1962}, \csource{find_points.c}{1976}{1979},
  \csource{find_points.c}{1982}{1985}), of a divisor (\csource{find_points.c}{1993}{2009}) and of
  the divisor shape (\csource{find_points.c}{1987}{1992}, \csource{find_points.c}{2010}{2016},
  \csource{find_points.c}{2019}{2023}), of the share that the checked loop keeps
  (\csource{find_points.c}{2051}{2065}) and of every denominator
  (\csource{find_points.c}{2025}{2027}, \csource{find_points.c}{2038}{2044},
  \csource{find_points.c}{2047}{2050}, \csource{find_points.c}{2066}{2066}); the estimate from the
  accumulators (\csource{find_points.c}{2074}{2105}) and its dispatch over the shapes
  (\csource{find_points.c}{1939}{1944}, \csource{find_points.c}{1959}{1959},
  \csource{find_points.c}{1986}{1986}, \csource{find_points.c}{2024}{2024},
  \csource{find_points.c}{2106}{2106}); and the mean bits per word of the bit arrays of the run, the
  number of set bits of a word a loop (\csource{find_points.c}{2128}{2132},
  \csource{find_points.c}{2143}{2151}, \csource{find_points.c}{2134}{2142},
  \csource{find_points.c}{2152}{2159}).
\end{definition}

The doubles of the cost model are Lean core's \verb|Float|, which has a logical model of its
arithmetic since Lean v4.35: \verb|+|, \verb|-|, \verb|*|, \verb|/| and \verb|sqrt| are the
correctly rounded IEEE~754 operations, as the C code compiled without fused multiply-adds computes
them, and the compiled Lean code, built for the baseline x86-64, fuses no multiply-adds either; the
conversions of natural numbers, integers and $64$-bit words to \verb|Float| (\verb|Float.ofNat|,
\verb|UInt64.toFloat|) round to nearest with ties to even, exact below $2^{53}$, as the conversions
of the C code do, and the one conversion of a double to a \verb|long|, in the offset of the second
stage, truncates a value in range by \verb|Float.toUInt64|; \verb|log|, \verb|pow|, \verb|ceil| and
\verb|floor| are the functions of the platform's C library, opaque to the logic. glibc's \verb|log|
and \verb|pow| are not correctly rounded, and glibc chooses their variant by the processor
(Section~\ref{sec:moduliC}), so that the choice depends on the machine and the model agrees with
ratpoints on one machine. On glibc the Lean runtime binds an older symbol version of \verb|log| and
\verb|pow| than ratpoints does, the two versions giving the same values on every argument that
occurs here ($x \ge 0$ for \verb|log|, a base in $(0, 1]$ and a positive exponent for \verb|pow|)
and differing only in \verb|errno| and in the sign of the NaN of the logarithm of a negative number;
$\log 2$, which gcc folds into the value of glibc at compile time, the model computes at run time,
with the same value; \verb|ceil| and \verb|floor|, exact, are calls of the library where gcc
compiles instructions. Every expression keeps the order of the operations of the C code, and every
comparison is the C code's, also where a NaN can arise (\verb|-W inf| with \verb|-A 2| makes
$0 \cdot \infty$ in the correction). The cost model has no contracts: no proof reads a value of a
\verb|Float|, so that every theorem holds for whatever values the model computes, and the choice is
free, for the proofs, among the values that satisfy the caps of the natural numbers of the C code. A
double that the caller gives (a tuning value, an end of a search interval) enters the arithmetic by
its bits (\lname{Ratpoints.Double.toFloat}, which joins Definition~\ref{def:Double} with I.3); a
double of the cost model that the report prints leaves it so (\lname{Ratpoints.Double.ofFloat},
which joins that node with I.15). No \verb|#guard| pins a value of \verb|log| or \verb|pow|, nor an
output that depends on the choice of the moduli, since the guards run on whatever machine builds the
project.

\section{The choice of the moduli}\label{sec:choice}

\begin{definition}[The options of the sieve]\label{def:sieveOptions}
  \uses{def:Double, def:choiceConstants}
  \lean{Ratpoints.SieveOptions, Ratpoints.SieveOptions.normalize, Ratpoints.SieveOptions.tuning}
  \leanok
  The fields of the arguments of the library that tune the
  sieve, as the caller gives them: the number of primes to look at, the cut points and the number of
  primes of the third stage, the bound of the forbidden divisors and the values of \verb|-R| and
  \verb|-A|, as integers for the C \verb|long|s, $-1$ when absent; the values of \verb|-r|,
  \verb|-U|, \verb|-C|, \verb|-Q| and \verb|-W|, doubles by their bits
  (Definition~\ref{def:Double}), $-1.0$ when absent (\lname{Ratpoints.Double.minusOne}, which joins
  that node with I.8); the array size, an integer, $\mathrm{defaultArraySize}$ by default, any value
  at most $0$ standing for it; and the Jacobi switch, a Boolean for the bit
  \verb|RATPOINTS_NO_JACOBI| (\csource{main.c}{337}{347}, \csource{main.c}{349}{350},
  \csource{main.c}{352}{352}, \csource{ratpoints.h}{365}{367}). Their normalization by the search:
  the number of primes $\mathrm{defaultNumPrimes}$ for a negative one and at most
  $\mathrm{numPrimes}$, the end of the second stage at most that number and the end of the first at
  most that of the second when the second is given, the bound of the forbidden divisors
  $\mathrm{defaultMaxForbidden}$ for a negative one and at most the size of the table
  (\csource{find_points.c}{4643}{4649}, \csource{find_points.c}{4655}{4656},
  \csource{find_points.c}{4662}{4665}), the array size normalized with the height bound where it is
  read ($\mathrm{arraySizeAt}$, Definition~\ref{def:arraySize}). The tuning of the choice: the
  target of the first stage when it is positive, the offset and the scale of the second and the cost
  of a row when they are at least $0$, as the C code tests them, and the default macro otherwise, a
  NaN giving the default (\csource{find_points.c}{2569}{2575}); the other tuning values are read
  where the C code reads them (Section~\ref{sec:moduliC}).
\end{definition}

The result of the choice, $\mathrm{Params}$ (Definition~\ref{def:Params}), becomes internal to the
search: the moduli in the order of the stages with the reserve last, how many primes the loop of the
tables examined ($\mathrm{examined}$, which bounds the first forbidden divisors and the factors of
the composite moduli) and how many were looked at in all ($\mathrm{lookedAt}$, which the report
prints), the cut points $s_1 \le s_2 \le s_3$, the bound of the forbidden divisors, the Jacobi
switch, the array size, and the estimates that the correction reads, the words and the denominators
of the run, the cost of a check relative to the reference and the words of the first correction, a
\verb|UInt64| (\csource{ratpoints.h}{368}{368}, \csource{ratpoints.h}{395}{397},
\csource{ratpoints.h}{406}{406}). The fields $s_3$, $\mathrm{lookedAt}$, the estimates, the cost of
a check and the words of the first correction come with the choice (I.8), and the parameters of the
executable fill them, until I.12 deletes those, with values that the run of the explicit route does
not read.

\begin{definition}[The composite moduli]\label{def:compositeModuli}
  \uses{def:costModel, def:densities, def:SieveEntry, def:mulMod, def:modinv, def:packingConstants}
  \lean{Ratpoints.primeEntry, Ratpoints.ModulusCandidate.ofPrime, Ratpoints.PowerTable, Ratpoints.powerIndex, Ratpoints.FacState, Ratpoints.takeFactor, Ratpoints.factorCandidate, Ratpoints.offerModulus, Ratpoints.addModuli, Ratpoints.makeModulus}
  \leanok
  The sieve entry of a prime, as \verb|examine_prime| makes it, a value with no C code of its own,
  and the candidate of a prime, its count $-1$ for the code $2$
  (\csource{find_points.c}{1782}{1785}). The prime powers examined so far, by the position of the
  prime and the exponent, with the records of the useful ones (\csource{find_points.c}{2350}{2350},
  \csource{find_points.c}{2352}{2354}, \csource{find_points.c}{2358}{2359}), and the index of a
  prime power, examined the first time (\csource{find_points.c}{2387}{2394}). The state of the
  factorization of a candidate (\csource{find_points.c}{2362}{2362},
  \csource{find_points.c}{2364}{2365}); a factor $p^e$ taken, refused as the factor too many or a
  prime beyond the first $64$, with the density of the prime from the loop of the tables or of the
  useful prime power (\csource{find_points.c}{2381}{2386}, \csource{find_points.c}{2395}{2398}); the
  factorization over the primes looked at, a cofactor below $p^2$ looked up among them
  (\csource{find_points.c}{2370}{2380}, \csource{find_points.c}{2399}{2399}); the candidate of an
  odd composite, none when its table is too dear, it does not factor, it is a prime or its density
  is not below $1$ (\csource{find_points.c}{2363}{2363}, \csource{find_points.c}{2368}{2368},
  \csource{find_points.c}{2400}{2402}, \csource{find_points.c}{2404}{2406}); the candidates offered,
  each keyed for the first stage, until the pool is full (\csource{find_points.c}{2346}{2349},
  \csource{find_points.c}{2351}{2351}, \csource{find_points.c}{2355}{2356},
  \csource{find_points.c}{2361}{2361}, \csource{find_points.c}{2403}{2403},
  \csource{find_points.c}{2407}{2411}); and the sieve entry of a composite modulus taken, from the
  codes of its factors: their tables, the product of their moduli, its reciprocal, bias and inverse
  of $\mathrm{rbaLength}$, the inverses of the strides empty, its reference and its builder
  (\csource{find_points.c}{2414}{2426}, \csource{find_points.c}{2428}{2429},
  \csource{find_points.c}{2434}{2470}, \csource{find_points.c}{2472}{2475}; the lines of the tables
  of the rows, 2427 and 2471, stay with Definition~\ref{def:storeTables}).
\end{definition}

\begin{definition}[The first two stages]\label{def:firstStages}
  \uses{def:costModel, def:compositeModuli, def:sieveOptions, def:runShape}
  \lean{Ratpoints.dropUsed, Ratpoints.takeEntries, Ratpoints.rankSecond, Ratpoints.dropComposites, Ratpoints.stageEntries, Ratpoints.firstStages}
  \leanok
  The candidates whose mask meets the primes taken dropped from a position on
  (\csource{find_points.c}{2509}{2518}); the candidates that one stage takes, in the order they are
  in, a given number or while the rule of the first stage wants the next, with the mask of the
  primes taken and the product of the densities (\csource{find_points.c}{2490}{2508},
  \csource{find_points.c}{2519}{2520}); the pool keyed again for the second stage and sorted from
  the end of the first (\csource{find_points.c}{2853}{2863}); the candidates after the second stage
  that are not primes saying more than coprimality dropped (\csource{find_points.c}{2879}{2890});
  the sieve entries of the composite moduli of the first two stages made
  (\csource{find_points.c}{2894}{2900}); and the first two stages with their cut points, the values
  of \verb|-n| and \verb|-N| when given (\csource{find_points.c}{2843}{2846},
  \csource{find_points.c}{2864}{2871}).
\end{definition}

\begin{definition}[The loop of the tables]\label{def:tablesLoop}
  \uses{def:densities, def:costModel, def:runShape, def:forbidden, def:storeTables,
    def:sieveOptions, def:compositeModuli, def:examinePrimeExit}
  \lean{Ratpoints.TableLoop, Ratpoints.examineInLoop, Ratpoints.raiseLimit, Ratpoints.tablesLoop, Ratpoints.choiceEstimate, Ratpoints.checkRelOf, Ratpoints.adaptStart}
  \leanok
  The state of the loop: the pool of the
  candidates, the number of primes to look at, the primes that say only coprimality, the table and
  the density of each prime that says something, whether denominators divisible by each prime occur,
  the forbidden divisors and the store (\csource{find_points.c}{2537}{2539},
  \csource{find_points.c}{2545}{2545}, \csource{find_points.c}{2547}{2547},
  \csource{find_points.c}{2566}{2567}); one prime examined in the loop, with the exit, its candidate
  keyed for the first stage when it says something and its forbidden divisor
  (\csource{find_points.c}{2606}{2626}, \csource{find_points.c}{2677}{2677}); the decision to look
  at one more prime, after the in-place sort of the pool by key
  (\csource{find_points.c}{2694}{2705}); the loop over the positions of the table that ends when the
  position reaches the number to look at, which the body raises, at most $\mathrm{numPrimes}$
  (\csource{find_points.c}{2601}{2605}, \csource{find_points.c}{2693}{2693},
  \csource{find_points.c}{2707}{2707}); what the choice derives from an estimate of the run, the
  cost of the calls per word, the bits per word swept and the offset of the second stage
  (\csource{find_points.c}{2597}{2599}, \csource{find_points.c}{2813}{2815}); the cost of a check
  relative to the reference, $1$ when it is not positive (\csource{find_points.c}{2587}{2588}); and
  the words of the first correction (\csource{find_points.c}{2580}{2581}).
\end{definition}

\begin{definition}[The third stage]\label{def:thirdStage}
  \uses{def:costModel, def:runShape, def:examinePrimeExit, def:compositeModuli, def:sieveOptions}
  \lean{Ratpoints.bestByR, Ratpoints.rTyp, Ratpoints.ThirdState, Ratpoints.lookFurther, Ratpoints.thirdStage, Ratpoints.survivorsAtThird, Ratpoints.thirdOf}
  \leanok
  The candidate of least density from a position
  on, the first on ties (\csource{find_points.c}{2955}{2961}); the mean density of the candidates
  that say more than coprimality, $1$ for none (\csource{find_points.c}{2982}{2991}); the state of
  the stage (\csource{find_points.c}{2922}{2922}, \csource{find_points.c}{2933}{2933}); one more
  prime looked at, with the exit when the curve has no points modulo it
  (\csource{find_points.c}{2976}{2978}, \csource{find_points.c}{2994}{3008}); the loop of the stage,
  the candidate taken moved to its end (\csource{find_points.c}{2949}{2950},
  \csource{find_points.c}{2973}{2975}, \csource{find_points.c}{2980}{2981},
  \csource{find_points.c}{2992}{2993}, \csource{find_points.c}{3009}{3016}); the survivors per
  denominator expected at its start (\csource{find_points.c}{2941}{2947}); and the stage after the
  first two, with the number of primes it wants and its costs per survivor and per denominator
  (\csource{find_points.c}{2926}{2932}).
\end{definition}

\begin{definition}[The choice of the moduli]\label{def:choice}
  \uses{def:tablesLoop, def:compositeModuli, def:firstStages, def:thirdStage, def:runShape,
    def:forbidden, def:Params, def:sieveOptions, def:storeTables}
  \lean{Ratpoints.SievingInfo, Ratpoints.choiceEnd, Ratpoints.choiceTail, Ratpoints.sievingInfo}
  \leanok
  What the choice gives the search: the
  result of the choice, the sieve entries of all its moduli, the reserve included, their
  reciprocals, the forbidden divisors tested with words and by their valuations, the store of the
  tables of rows and the switch of the checked loop (\csource{ratpoints.h}{389}{389},
  \csource{ratpoints.h}{391}{391}); the end of the choice: the third stage and the reserve sorted by
  density from the end of the second stage on, the entries of all the candidates in their order,
  their reciprocals and the store with the tables of all of them
  (\csource{find_points.c}{3017}{3017}, \csource{find_points.c}{3024}{3034},
  \csource{find_points.c}{3064}{3070}); the choice after the loop of the tables: the estimate again
  with the forbidden divisors, the keys again, the composite moduli offered, the sort by key, the
  first two stages, the entries and the third stage (\csource{find_points.c}{2808}{2812},
  \csource{find_points.c}{2816}{2823}, \csource{find_points.c}{2826}{2826}); and
  \verb|sieving_info|: the tuning, the first estimate, the loop of the tables, the further forbidden
  divisors, the switch of the checked loop cleared when nothing is left for it to test, and the
  rest, with the exit prime as an error (\csource{find_points.c}{2522}{2532},
  \csource{find_points.c}{2546}{2546}, \csource{find_points.c}{2548}{2559},
  \csource{find_points.c}{2595}{2596}, \csource{find_points.c}{2794}{2795},
  \csource{find_points.c}{3102}{3103}).
\end{definition}

The function keeps the name of the explicit $\mathrm{sievingInfo}$ of stage~F1
(Definition~\ref{def:sievingInfo}), which made the entries of given moduli, and its contract the
names of the contract of that one (\verb|IsSievingInfo|, \verb|sievingInfo_spec|,
Lemma~\ref{lem:sievingInfo}): I.8 renames the explicit function and its contract to
\lname{Ratpoints.sievingInfoOf}, \lname{Ratpoints.IsSievingInfoOf} and
\lname{Ratpoints.sievingInfoOf\_spec}, which then no longer model C code and which
Definition~\ref{def:sievingInfo} and Lemma~\ref{lem:sievingInfo} carry until I.12 deletes them with
the explicit route. The C ranges move to the components above from the definitions of the explicit
route, each with the item that brings its new owner. With I.3 the lines 1619-1623 of
\verb|find_points.c| go from the parameters of the executable (Definition~\ref{def:sieveParams}) to
the offset of the second stage (Definition~\ref{def:costModel}). With I.7: from the parameters of
the executable the lines 2503, 2845-2846, 2864-2869, 2922, 2949, 2973-2975, 2980 and 3015; from the
factors of a modulus (Definition~\ref{def:modulusFactors}) its lines 2362, 2365, 2370-2386 and
2396-2400; from the sieve entry of a modulus (Definition~\ref{def:makeModulus}) its lines 2350-2354,
2358-2359, 2387-2395, 2414-2426, 2428-2429, 2434-2470, 2472-2475, 2976-2978 and 2994-2998; and from
the explicit $\mathrm{sievingInfo}$ its lines 2894-2900. With I.8: from the parameters of the
executable the lines 2566, 2571-2572, 2599, 2693-2705, 2808, 2815, 2926-2928 and 3017, and
4643-4649, 4655-4656, 4662-4665 and \verb|main.c:347| to Definition~\ref{def:sieveOptions}; from the
loop of the tables of stage~F1 (Definition~\ref{def:examinePrimes}) its lines 2537, 2545, 2601-2602,
2605-2614, 2616, 2619, 2622, 2624-2626, 2677 and 2707; from the sieve entry of a modulus its lines
2546 and 2548; from the explicit $\mathrm{sievingInfo}$ its lines 2522-2532, 3064-3070 and
3102-3103; and from the search of the polynomial (Definition~\ref{def:searchPoints}) its lines
2794-2795, the switch of the checked loop. The line 4577 and \verb|main.c:221-222| go with I.10 to
the search and to the program. The comment lines among them (2369, 2793, 2892-2893) are left to no
declaration. Nothing of the explicit route models C code after I.10.

\begin{lemma}[The pool of the candidates]\label{lem:choicePool}
  \uses{def:costModel, def:firstStages, def:thirdStage, def:sieveOptions}
  \lean{Ratpoints.AllQ, Ratpoints.AllQ.push, Ratpoints.AllQ.eraseIdxIfInBounds, Ratpoints.AllQ.sortFrom, Ratpoints.AllQ.mapIdx_key, Ratpoints.AllQ.rankSecond, Ratpoints.TailNf, Ratpoints.DropUsedInvariant, Ratpoints.TakeEntriesInvariant, Ratpoints.DropCompositesInvariant, Ratpoints.dropUsed_spec, Ratpoints.takeEntries_spec, Ratpoints.dropComposites_spec, Ratpoints.bestByR_spec, Ratpoints.firstStages_spec, Ratpoints.firstStages_run, Ratpoints.size_sortFrom, Ratpoints.getElem_sortFrom_of_lt, Ratpoints.exists_getElem_sortFrom_eq, Ratpoints.size_rankSecond, Ratpoints.map_some_filterMap_ssp, Ratpoints.length_filterMap_ssp_and_forall, Array.Extends, Array.Extends.refl, Array.Extends.trans, Array.extends_push, Array.extends_ite_push, Array.getElem!_eq_of_getElem?_eq_some, Option.eq_none_of_forall_prod4}
  \leanok
  For a property $Q$ of the candidates that does not read their key and cost, every candidate of the
  pool has $Q$ after a push, a deletion, a sort and a keying again; the loops of the stages keep it,
  the first stage ends at most at the size of the pool and the second after it; after the deletion
  of the composite moduli not taken every candidate from the end of the second stage on counts no
  factor; the candidate of least density lies from the end of the third stage on; and the first two
  stages give $s_1 \le s_2 \le$ the size of the pool, whatever the keys are. A sort keeps the size
  and moves the candidates only among themselves.
\end{lemma}

\begin{proof}
  \leanok
  The loops of the stages delete and keep candidates only; their invariants are that every candidate
  has $Q$ and the position is within the pool, with the measures the size of the pool less the
  position. The sorts are permutations (Lean's merge sort), which keep every property of the
  members.
\end{proof}

\begin{lemma}[The composite moduli are sound]\label{lem:compositeModuli}
  \uses{def:compositeModuli, def:thirdStage, lem:densities, lem:choicePool, def:IsPoint}
  \lean{Ratpoints.factorOfCode, Ratpoints.PrimeSeOK, Ratpoints.PowerOK, Ratpoints.CodeOK, Ratpoints.CodeOK.isSound, Ratpoints.CodeOK.extend, Ratpoints.CodeOK.factorOfCode_extend, Ratpoints.ModulusCandidate.IsPrime, Ratpoints.ModulusCandidate.IsComposite, Ratpoints.ModulusCandidate.IsComposite.isSound, Ratpoints.CandOK, Ratpoints.CandOK.of_eq, Ratpoints.isPrime_ofPrime, Ratpoints.PowerTable.OK, Ratpoints.FacState.OK, Ratpoints.powerIndex_spec, Ratpoints.takeFactor_spec, Ratpoints.factorCandidate_spec, Ratpoints.offerModulus_spec, Ratpoints.AddModuliInvariant, Ratpoints.addModuli_spec, Ratpoints.EntryOK, Ratpoints.entryOK_primeEntry, Ratpoints.makeModulus_spec, Ratpoints.ChoicePoolOK, Ratpoints.StageEntriesInvariant, Ratpoints.stageEntries_spec, Ratpoints.ThirdState.OK, Ratpoints.ThirdStageInvariant, Ratpoints.lookFurther_spec, Ratpoints.thirdStage_spec}
  \leanok
  A code of a factor is valid when it is the position of a prime of the table whose table the loop
  of the tables has, with its entry, or $\mathrm{numPrimes}$ plus the index of a useful prime power
  examined with its records; the factor that it names is then sound for every modulus
  $0 < m < \mathrm{maxPrimeEven}$ that the modulus of the factor divides. A candidate is a prime of
  the table at its position with its prime entry, or a composite: odd,
  $1 < m < \mathrm{maxPrimeEven}$, with valid codes whose moduli multiply to $m$, no entry and as
  many factors counted as codes; this does not depend on the key, the cost, the density or the mask.
  The factorization, the offer and the loop of the composite moduli keep the pool of such
  candidates, and the table of the prime powers only grows, which keeps the codes valid; the entry
  of a composite modulus made from its codes is a sound modulus entry, with the reference of its
  position and a prime entry when it counts no factor; the entries of the first two stages are so
  made; and the third stage keeps the pool so, every candidate from the end of the second stage on a
  prime, and its exit at a further prime means that the curve has no points.
\end{lemma}

\begin{proof}
  \uses{lem:examinePrime_eq, lem:examinePrime_none, def:stage3, lem:stage3}
  \leanok
  The contracts follow the loops of \verb|add_moduli| (three nested, with the measures
  $2(\mathrm{pnLim} - \mathrm{pn}) + \mathrm{ok}$, $\mathrm{pnLim} - \mathrm{pn}$ and the cofactor)
  and of the third stage (with the measure
  $(\mathrm{numPrimes} - \mathrm{pnLim}) + (\mathrm{want} - s_3)$). The soundness of a composite
  entry is that of Lemma~\ref{lem:densities}, for the factors named by the codes, whose moduli
  divide $m$ and multiply to it; the soundness of a factor reads only the divisibility. The exit of
  the third stage is that of the examination of a prime (Lemma~\ref{lem:examinePrime_none}).
\end{proof}

\begin{lemma}[The choice is what the run needs]\label{lem:choice}
  \uses{def:choice, lem:choicePool, lem:compositeModuli, lem:densities, lem:forbidden,
    lem:examinePrime_none, lem:storeTables, def:decisionWF}
  \lean{Ratpoints.IsSievingInfo, Ratpoints.sievingInfo_spec, Ratpoints.TableLoop.OK, Ratpoints.TableLoop.ok_init, Ratpoints.TablesLoopInvariant, Ratpoints.examineInLoop_spec, Ratpoints.raiseLimit_spec, Ratpoints.tablesLoop_spec, Ratpoints.EntriesOK, Ratpoints.entriesOK_choiceEnd, Ratpoints.isSievingInfo_choiceEnd, Ratpoints.ChoicePoolOK.sortFrom, Ratpoints.ChoicePoolOK.exists_entryOK, Ratpoints.firstStages_post, Ratpoints.thirdOf_spec, Ratpoints.choiceTail_spec, Ratpoints.Sieve.RowStore.MadeFor.foldl_addTables, Ratpoints.SieveOptions.numPrimes_normalize_le, Ratpoints.SieveOptions.maxForbidden_normalize_nonneg}
  \leanok
  For options that look at no more
  than $\mathrm{numPrimes}$ primes and test a nonnegative number of forbidden divisors, as the
  normalized options do, the checked loop only for even degree, and any values of the doubles: an
  exit of the choice, at a prime of the loop of the tables or of the third stage without points
  modulo it, means that the curve has no points of any height; and otherwise the result is what the
  run needs ($\mathrm{IsSievingInfo}$): every sieve entry sound, with the fields that the packing
  reads and the tables of inverses of its factors; an entry that counts no factor a prime entry of
  the table; the reciprocals those of the moduli of the entries, in their order; the first decision
  $(0, s_2, s_3)$ well formed (Definition~\ref{def:decisionWF}: $s_1 \le s_2 \le s_3 \le$ the number
  of the entries, and every entry from $s_2$ on counting no factor); the forbidden divisors sound;
  the store with the tables of all the entries, the reserve included; and the checked loop only for
  even degree.
\end{lemma}

\begin{proof}
  \leanok
  The loop of the tables keeps at most $\mathrm{numPrimes}$ primes to look at, the tables of the
  primes those of the reference at their positions, every candidate a prime, the forbidden divisors
  sound and the store that of the primes examined (the invariant of the loop,
  \lname{Ratpoints.TablesLoopInvariant}, that of the loop of stage~F1 with the candidates); the
  composite moduli and the stages keep the pool of Lemma~\ref{lem:compositeModuli}, and at the end
  of the choice every candidate has its entry, so that the entries line up with the candidates and
  the cut points with both. The doubles enter only through values that no conjunct reads.
\end{proof}

The proofs need no well-formedness of $\mathrm{Params}$ any more: $\mathrm{Params.WF}$, a hypothesis
of every theorem since stage~C1, is replaced by $\mathrm{IsSievingInfo}$ for the result of the
choice and by the well-formedness of each decision of the correction (Lemma~\ref{lem:correctionWF}),
so that the route of the proofs is the contract of the run for every result of the choice that has
$\mathrm{IsSievingInfo}$ and every decision that is well formed. I.12 deletes $\mathrm{Params.WF}$,
its decidability instance and its lemmas, with \lname{Ratpoints.getElem!\_primeTable\_inj}, which no
proof uses then, and rewords the docstrings of the model that name it.

\section{The counters}\label{sec:counters}

\begin{definition}[The counters of the sieve]\label{def:counters}
  \lean{Ratpoints.Sieve.Counters}
  \leanok
  The six counters of the sieving thread, the words swept by the first stage, the bit arrays that
  survive it, the numerators extracted, those coprime to the denominator and those that pass the
  third stage, and the calls of the sieve of one denominator, as $64$-bit words that wrap as the
  \verb|unsigned long|s of the C code do (\csource{rp-private.h}{579}{581}): a record of its own,
  passed beside the delivery and returned with it, so that the delivery of stage~G, and with it the
  initial delivery $\langle \mathrm{info}, 0, \mathrm{false}\rangle$ of the protocol statement, does
  not change. The components count where the C code counts, within the lines that they model
  already: the test of a candidate (Definition~\ref{def:accepted}) the coprime numerators after the
  test for common factors and the checks after the third stage, returning the counters with its
  verdict; the walk through the set bits of a word (Definition~\ref{def:walkWord}) each numerator
  extracted, in both arms of the extraction (Definition~\ref{def:siftBitArray}); the scan of a piece
  (Definition~\ref{def:scanPiece}) each bit array that is not zero after the first stage; one piece
  (Definition~\ref{def:siftPiece}) its $n\,\mathrm{rbaPack}$ words, multiplied as \verb|UInt64|s as
  the C code multiplies the \verb|unsigned long| by \verb|RBA_PACK|; the bits of an interval
  (Definition~\ref{def:siftBits}) passes them through its pieces; and one denominator
  (Definition~\ref{def:sift}) its call, gaining the link \verb|rp-private.h:585|, the fields of the
  sieving thread that it keeps.
\end{definition}

\begin{lemma}[The counters in closed form]\label{lem:counters}
  \uses{def:counters, lem:accepted, lem:siftBitArray, lem:relprime, lem:stage3, def:exactBounds,
    def:scanPiece, def:siftBits, def:siftPiece, lem:reads_row}
  \lean{Ratpoints.Sieve.instAddCounters, Ratpoints.Sieve.Counters.ext, Ratpoints.Sieve.Counters.add_assoc, Ratpoints.Sieve.Counters.add_comm, Ratpoints.Sieve.Counters.add_empty, Ratpoints.Sieve.Counters.empty_add, Ratpoints.Sieve.testedCounts, Ratpoints.Sieve.testedCounts_nil, Ratpoints.Sieve.testedCounts_append, Ratpoints.Sieve.testedCounts_eq, Ratpoints.Sieve.snd_run_accepted, Ratpoints.Sieve.run_relprime_toInt64, Ratpoints.Sieve.walkWord_counters_spec, Ratpoints.Sieve.siftBitArrayByWords_counters_spec, Ratpoints.Sieve.siftBitArrayWhole_counters_spec, Ratpoints.Sieve.siftBitArray_counters_spec, Ratpoints.Sieve.isZeroAt_eq_not_any, Ratpoints.Sieve.scanPiece_counters_spec, Ratpoints.Sieve.arraysWith, Ratpoints.Sieve.siftPiece_counters_spec, Ratpoints.Sieve.bitsCounts, Ratpoints.Sieve.siftBits_counters_spec, Ratpoints.Sieve.intervalCounts, Ratpoints.Sieve.siftCounts, Ratpoints.Sieve.SiftIntervalsCounts, Ratpoints.Sieve.WalkCountsInvariant, Ratpoints.Sieve.SiftBitArrayCounts, Ratpoints.Sieve.ScanPieceCounts, Ratpoints.Sieve.SiftCounts}
  \leanok
  The counters add field by field, wrapping, a commutative and associative operation with the empty
  record as zero. While the callback has not stopped the search, each counting component returns its
  counters plus a function of its arguments that reads neither the callback, nor the delivery, nor
  the flags of the output, nor the store, nor the size of the pieces: for the test of a list of
  numerators of a denominator $b$, their number, those coprime to $b$ and those that pass the test
  for common factors and the third stage, which for $|a| \le H < 2^{63}$, $0 < b < 2^{63}$ and prime
  entries of the third stage are the numerators with $\gcd(a, b) = 1$ and those of them that survive
  the entries; for a bit array, the test of the numerators of its set bits; for a piece, its words,
  the bit arrays that hold a bit that passes the class test, the first stage and the masks, and the
  test of the numerators that pass all the stages; for the bits of an interval, the words of all its
  bit arrays, the bit arrays holding a bit in the interval that passes the first stage, and the test
  of the numerators in the interval that pass all the stages; and for one denominator, one call plus
  the counts of its intervals up to the first that ends the denominator. The contract of one
  denominator (Definition~\ref{def:sift}) gains this as a third conjunct, beside the store and the
  delivery.
\end{lemma}

\begin{proof}
  \uses{def:andRowsChunked, def:intervalBits, lem:andRowsChunked}
  \leanok
  Each counting component has a contract of its own for its counters, with named invariants that
  carry the counts of the part of the loop done, so that the contracts of the delivery keep their
  proofs; the contract of one denominator passes those of its callees to the program logic. The
  counts are counted before the exact check, from the same bits that the contracts of the delivery
  read, so that they are lengths of the lists that those contracts already filter. No contract of
  the counters needs the callback not to have stopped before, only that it does not stop within.
\end{proof}

\section{The blocks}\label{sec:blocks}

\begin{definition}[The blocks of a run]\label{def:blockPlan}
  \uses{def:Shape, def:choiceConstants}
  \lean{Ratpoints.Sieve.Shape.segments, Ratpoints.Sieve.BlockPlan, Ratpoints.Sieve.BlockPlan.nblocks, Ratpoints.Sieve.blockLength, Ratpoints.Sieve.Shape.blockPlan, Ratpoints.Sieve.BlockPlan.steps}
  \leanok
  The segments of steps of a shape from $b_{\mathrm{low}}$ to $b_{\mathrm{high}}$, each from its
  first step to its last, and the number of steps of all of them: the $k$ with
  $b_{\mathrm{low}} \le k^2 \le b_{\mathrm{high}}$ for the squares, for each divisor $s$ of the
  divisor shape in turn a segment of the $k$ with
  $b_{\mathrm{low}} \le s k^2 \le b_{\mathrm{high}}$, the words from
  $\lfloor b_{\mathrm{low}}/64 \rfloor$ to $\lfloor b_{\mathrm{high}}/64 \rfloor$, and the
  denominators themselves (\csource{find_points.c}{3812}{3837}); the blocks of a run, its segments,
  the length of its blocks and the first block of each segment, with the number of the blocks at the
  end (\csource{find_points.c}{3750}{3754}); the length of the blocks,
  $\lfloor \mathrm{steps}/\mathrm{blocksWanted} \rfloor + 1$, at least $\mathrm{blockMin}$, and the
  fixed length when the build sets one (\csource{find_points.c}{3838}{3842}); the plan of the
  blocks, a segment with steps holding
  $\lfloor (\mathrm{hi} - \mathrm{lo})/\mathrm{len} \rfloor + 1$ blocks
  (\csource{find_points.c}{3843}{3850}); and the segment and the first and last steps of block $i$,
  the first segment that ends after block $i$, the block of $\mathrm{len}$ steps or up to the end of
  its segment (\csource{find_points.c}{3900}{3904}).
\end{definition}

\begin{definition}[The loops of the shapes over a block]\label{def:blockLoops}
  \uses{def:blockPlan, def:sift, def:fillBpList, def:Shape, def:mulMod, def:useCLong, def:ctz,
    def:counters, def:forbidden}
  \lean{Ratpoints.Sieve.BlockState, Ratpoints.Sieve.siftDen, Ratpoints.Sieve.squaresBlock, Ratpoints.Sieve.squares1Block, Ratpoints.Sieve.forbStarts, Ratpoints.Sieve.denWord, Ratpoints.Sieve.walkDens, Ratpoints.Sieve.wordsBlock, Ratpoints.Sieve.allBlock, Ratpoints.Sieve.Shape.blockLoop}
  \leanok
  What a block carries from one denominator to the next, the delivery, the store and the counters;
  one denominator of a block, the residues of the first two stages (Definition~\ref{def:fillBpList})
  and its points (Definition~\ref{def:sift}), with the entries and the reciprocals of the decision
  of the block (\csource{find_points.c}{3926}{3927}, \csource{find_points.c}{3955}{3956},
  \csource{find_points.c}{4042}{4043}, \csource{find_points.c}{4066}{4067}); the steps of a block of
  the squares, of a segment of the divisor shape with its valuation test, and of every denominator,
  each denominator with a pattern handed to the step, returning when the callback has stopped the
  search (\csource{find_points.c}{3919}{3925}, \csource{find_points.c}{3928}{3935};
  \csource{find_points.c}{3937}{3943}, \csource{find_points.c}{3954}{3954},
  \csource{find_points.c}{3957}{3965}; \csource{find_points.c}{4060}{4065},
  \csource{find_points.c}{4068}{4074}); in the shape of the words, the positions of the forbidden
  divisors at the first word $a$ of the block, $a \bmod p$ by the reciprocal for
  $a \le \mathrm{mulModLimit}$ and by \verb|mod| otherwise (\csource{find_points.c}{3981}{3981},
  \csource{find_points.c}{3983}{3988}), the word of the denominators of a word that pass the class
  test and those divisors, the positions stepping on and the first and the last word of the run
  masked (\csource{find_points.c}{3990}{3990}, \csource{find_points.c}{3993}{4006}), the walk
  through its set bits with the tests of the shape (\csource{find_points.c}{4007}{4009},
  \csource{find_points.c}{4012}{4012}, \csource{find_points.c}{4044}{4053}) and the loop over the
  words (\csource{find_points.c}{3967}{3980}, \csource{find_points.c}{3982}{3982},
  \csource{find_points.c}{3989}{3989}, \csource{find_points.c}{3991}{3992},
  \csource{find_points.c}{4054}{4058}); and the dispatch over the shapes, the step an argument
  (\csource{find_points.c}{3917}{3918}, \csource{find_points.c}{3936}{3936},
  \csource{find_points.c}{3966}{3966}, \csource{find_points.c}{4059}{4059},
  \csource{find_points.c}{4075}{4076}).
\end{definition}

\begin{definition}[One block]\label{def:runBlock}
  \uses{def:blockLoops, def:counters}
  \lean{Ratpoints.Sieve.Decision, Ratpoints.Sieve.findDecision, Ratpoints.Sieve.BlockStat, Ratpoints.Sieve.runBlock}
  \leanok
  A decision of the correction, the first block it is for and the ends $s_2$ and $s_3$ from it on,
  an array of them for the three arrays of the C code; the decision in force for block $i$, the last
  whose first block is at most $i$, found from the last back (\csource{find_points.c}{3908}{3910},
  \csource{find_points.c}{3912}{3912}); the report of a block, its counters with the index and the
  ends $s_2$ and $s_3$ of its decision (\csource{find_points.c}{4082}{4086}); and block $i$: its
  segment and steps, the decision in force, the entries and the reciprocals of the first two and of
  the third stage cut by the ends $s_2$ and $s_3$ of that decision, the counters from zero, the loop
  of its shape and its report (\csource{find_points.c}{3127}{3128},
  \csource{find_points.c}{3889}{3898}, \csource{find_points.c}{3911}{3911},
  \csource{find_points.c}{3913}{3915}, \csource{find_points.c}{4077}{4078},
  \csource{find_points.c}{4088}{4093}, \csource{rp-private.h}{582}{582}).
\end{definition}

\begin{definition}[A well-formed decision]\label{def:decisionWF}
  \uses{def:runBlock, def:SieveEntry}
  \lean{Ratpoints.Sieve.DecisionWF, Ratpoints.Sieve.DecisionWF.mk_sp2, Ratpoints.Sieve.DecisionWF.map_dinv}
  \leanok
  A decision is well formed for the sieve entries $E$ and the end $s_1$ of the first stage when
  $s_1 \le s_2 \le s_3 \le |E|$ and every entry from $s_2$ on counts no factor, so that it is a
  prime entry of the table for a result of the choice. A well-formed decision stays so with
  $s_3 = s_2$ and for the entries with other inverses of the strides.
\end{definition}

\begin{lemma}[The plan of the blocks]\label{lem:blockPlan}
  \uses{def:blockPlan}
  \lean{Ratpoints.Sieve.numBlocks, Ratpoints.Sieve.blockRange, Ratpoints.Sieve.blockRange_succ, Ratpoints.Sieve.blockRange_le, Ratpoints.Sieve.numBlocks_of_le, Ratpoints.Sieve.flatMap_blockRange, Ratpoints.Sieve.Shape.numSegs, Ratpoints.Sieve.Shape.segment, Ratpoints.Sieve.Shape.segments_spec, Ratpoints.Sieve.Shape.toList_fst_run_segments, Ratpoints.Sieve.BlockPlan.WF, Ratpoints.Sieve.Shape.blockPlan_spec, Ratpoints.Sieve.BlockPlan.WF.nblocks_eq, Ratpoints.Sieve.BlockPlan.WF.segBlock_mono, Ratpoints.Sieve.SegmentSearchInvariant, Ratpoints.Sieve.BlockPlan.WF.steps_spec, Ratpoints.Sieve.BlockPlan.WF.run_steps, Ratpoints.Sieve.BlockPlan.WF.run_steps_eq, Ratpoints.Sieve.BlockPlan.WF.run_steps_le, Ratpoints.Sieve.BlockPlan.blockSteps, Ratpoints.Sieve.BlockPlan.blockSeg, Ratpoints.Sieve.BlockPlan.WF.flatMap_blockSteps, Ratpoints.Sieve.one_le_blockLength}
  \leanok
  The length of the blocks is at least $1$; the blocks of a segment from $\mathrm{lo}$ to
  $\mathrm{hi}$ have steps within it and, in order, cover its steps in order; the plan of a shape is
  well formed, its first blocks the sums of the numbers of the blocks of the segments before,
  increasing, with the number of the blocks at the end; for $i < \mathrm{nblocks}$ the steps of
  block $i$ are those of its segment by the blocks of that segment, with $a \le e$, which the loops
  over the steps need to end; and the blocks in order, each mapped by a function of its segment,
  give the steps of the segments in order.
\end{lemma}

\begin{proof}
  \leanok
  The search for the segment of a block ends by the measure of the segments left; the rest is
  arithmetic of the quotients by the length, which is positive by the decided fact on the fixed
  length.
\end{proof}

\begin{lemma}[The loops of the shapes]\label{lem:blockLoops}
  \uses{def:blockLoops, lem:mulMod, lem:forbidden}
  \lean{Ratpoints.Sieve.BlockPost, Ratpoints.Sieve.BlockInv, Ratpoints.Sieve.MapBlockInvariant, Ratpoints.Sieve.AllBlockInvariant, Ratpoints.Sieve.WalkDensWhileInvariant, Ratpoints.Sieve.WordsBlockInvariant, Ratpoints.Sieve.squaresBlock_spec, Ratpoints.Sieve.squares1Block_spec, Ratpoints.Sieve.allBlock_spec, Ratpoints.Sieve.walkDens_spec, Ratpoints.Sieve.ForbFits, Ratpoints.Sieve.PosAt, Ratpoints.Sieve.posAt_forbStarts, Ratpoints.Sieve.denWord_spec, Ratpoints.Sieve.wordDens, Ratpoints.Sieve.wordBits_eq_foldl, Ratpoints.Sieve.wordsBlock_spec, Ratpoints.Sieve.Shape.stepDens, Ratpoints.Sieve.Shape.ForbFits, Ratpoints.Sieve.Shape.LoopFits, Ratpoints.Sieve.Shape.blockLoop_spec}
  \leanok
  For a step with a contract on each denominator with a given property (for every member of a family
  of lists it delivers that member's list of the denominator, keeps a property of the store and adds
  the counts of the denominator), the loop of each shape over the steps $a \le e$ of a block, when
  the denominators of those steps that pass the test of the shape have the property, delivers the
  concatenation of the lists of these denominators, in order, keeps the property and adds the counts
  in the same order, returning after a step at which the callback stopped the search. In the shape
  of the words, for forbidden divisors with $0 < p \le 2^{32}$ and their reciprocals, the positions
  at the first word of the block are $a \bmod p$, and the word of the denominators of the word $w$
  that the class test and those divisors leave, with the first and the last word of the run masked,
  is that of the listing of stage~E, and the positions are then those at $w + 1$, so that the words
  of a block give the denominators of the listing in those words.
\end{lemma}

\begin{proof}
  \uses{lem:counters}
  \leanok
  The invariants say that the delivery is that of the prefix of the steps done and the counts their
  sum, for every member of the family, which need not have one, and the property of the store; the
  stop of the callback ends the loop with the delivery of a prefix, which is that of the whole list
  by the delivery after a stop (Lemma~\ref{lem:deliver}). The positions by the reciprocal are those
  of the reduction (Lemma~\ref{lem:mulMod}).
\end{proof}

\begin{lemma}[The blocks list the denominators of the shape]\label{lem:blockDens}
  \uses{def:blockPlan, lem:blockPlan, lem:blockLoops, def:Shape}
  \lean{Ratpoints.Sieve.blockDens, Ratpoints.Sieve.Shape.flatMap_blockDens, Ratpoints.Sieve.Shape.flatMap_range_numSegs, Ratpoints.Sieve.Shape.mem_dens_of_mem_blockDens, Ratpoints.Sieve.Shape.flatMap_filter_blockDens, Ratpoints.Sieve.Shape.eq_of_mem_blockDens}
  \leanok
  The denominators of the steps of block $i$, before the tests of the shape, concatenated over the
  blocks of the run in order, are the denominators of the shape from $b_{\mathrm{low}}$ to
  $b_{\mathrm{high}}$ in the order of its loop, the listing of stage~E (Definition~\ref{def:Shape}),
  a definition of the model until the run sifts the blocks (I.11) and of the proofs from then on;
  lists made block by block that agree on each block with lists of the denominators give those lists
  over the listing filtered by the test of the shape; and when the listing has no duplicates, a
  denominator lies in at most one block.
\end{lemma}

\begin{proof}
  \leanok
  By Lemma~\ref{lem:blockPlan}, the blocks in order give the steps of the segments in order; mapped
  by the denominators of the steps of each segment, these are the listing of each shape, segment by
  segment for the divisor shape, word by word for the shape of the words by
  Lemma~\ref{lem:blockLoops}.
\end{proof}

\begin{lemma}[One block]\label{lem:runBlock}
  \uses{def:runBlock, def:decisionWF, lem:blockLoops, lem:blockDens, lem:counters, lem:buildRows,
    lem:fetchRow, lem:forbidden, lem:fillBpList}
  \lean{Ratpoints.Sieve.RunWords, Ratpoints.Sieve.denPoints, Ratpoints.Sieve.denCounts, Ratpoints.Sieve.take_append_take_drop, Ratpoints.Sieve.siftDen_spec, Ratpoints.Sieve.runBlock_spec, Ratpoints.Sieve.decisionWF_getElem!_findDecision, Ratpoints.Sieve.IsForbiddenEntry.mem_primeTable, Ratpoints.Sieve.forbFits_of_forall_isForbiddenEntry, Ratpoints.Sieve.forbFits_fst_run_runSetup, Ratpoints.Sieve.Shape.loopFits_run_blockPlan, Ratpoints.Sieve.RowStore.Keeps, Ratpoints.Sieve.RowStore.Keeps.refl, Ratpoints.Sieve.RowStore.Keeps.trans, Ratpoints.Sieve.RowStore.keeps_of_fetchRow, Ratpoints.Sieve.buildRows_keeps, Ratpoints.Sieve.StoreHolds.buildRows_setRows_take, Ratpoints.Sieve.RowStore.Fits.fits_run_sift, Ratpoints.Sieve.sift_storeHolds, Ratpoints.Sieve.DecisionSearchInvariant, Ratpoints.Sieve.findDecision_spec}
  \leanok
  Building rows never removes a row, so that a store that fits all the entries, the reserve
  included, still fits them after any denominator, whose sieve reads and builds the rows of its
  first two stages only. For a callback that has not stopped the search, an array size $A \ge 1$,
  sieve entries whose factors have their tables of inverses and whose entries that count no factor
  are prime entries, the reciprocals of their moduli, the row offsets of the classes for all the
  entries, the word of the classes that admit a numerator, a well-formed decision in force for block
  $i$ among the decisions, the plan of the shape, forbidden divisors with their reciprocals,
  $i < \mathrm{nblocks}$, $1 \le b_{\mathrm{low}}$, $b_{\mathrm{high}} < 2^{63}$, $H < 2^{63}$ and a
  store that fits all the entries: block $i$ keeps the store so, reports the decision in force with
  its ends $s_2$ and $s_3$, delivers the points of the denominators of its steps that pass the test
  of the shape, in order, each with the first $s_3$ entries of that decision, for every choice of
  the words of the classes and the ranges of the numerators that the packing and the bounds give,
  and, while the callback does not stop the search, reports the sum of the counts of those
  denominators with the ends $s_2$ and $s_3$ of the decision. The decision in force for block $i$,
  among decisions of which the first is for block $0$, is one whose first block is at most $i$ while
  every later one is for a block after $i$, and it is well formed when every decision is.
\end{lemma}

\begin{proof}
  \uses{lem:densities, lem:shape_dens}
  \leanok
  The step is one denominator, whose contract is that of Definition~\ref{def:sift} with the entries
  of the first two stages and of the third cut by the decision, $E_{:s_2}$ and
  $(E_{s_2:})_{:s_3 - s_2}$, which concatenate to $E_{:s_3}$, and with the frame of the store; the
  loop of the shape (Lemma~\ref{lem:blockLoops}) does the rest, and the block is straight-line code
  around it. The forbidden divisors of the choice are primes of the table, below $2^{32}$, with
  their reciprocals (Lemma~\ref{lem:densities}).
\end{proof}

\section{The correction of the moduli and the run}\label{sec:correction}

\begin{definition}[The correction of the moduli]\label{def:correction}
  \uses{def:runBlock, def:counters, def:costModel, def:sieveOptions, def:Params}
  \lean{Ratpoints.Sieve.RunCounts, Ratpoints.Sieve.RunCounts.fold, Ratpoints.Sieve.adaptRates, Ratpoints.Sieve.adaptUp, Ratpoints.Sieve.adaptDown, Ratpoints.Sieve.adaptSecond, Ratpoints.Sieve.adaptThird, Ratpoints.Sieve.adaptPrimes, Ratpoints.Sieve.Correction, Ratpoints.Sieve.Correction.start, Ratpoints.Sieve.runFold}
  \leanok
  The counters of the run, with the words swept when those from the extracted numerators on were
  last restarted and the end $s_2$ that they were counted under
  (\csource{find_points.c}{1363}{1367}), and their fold with a block, restarting those counters when
  the block was sifted with another $s_2$ (\csource{find_points.c}{4107}{4113}); the rates that the
  run has shown, the bit arrays per word after the first stage and the numerators per word after the
  second with the products of the densities of the stages, none when there is too little data
  (\csource{find_points.c}{1389}{1389}, \csource{find_points.c}{1400}{1420}); the walk of the end of
  the second stage up through the further entries while the survivors that the next removes are
  worth its cost (\csource{find_points.c}{1449}{1457}), and down through the entries below it while
  the position is above $s_1$, the entry counts no factor and dropping it costs less than it saves
  (\csource{find_points.c}{1460}{1472}); the correction of the second stage, which walks up and,
  when that takes nothing, down (\csource{find_points.c}{1439}{1448},
  \csource{find_points.c}{1458}{1459}, \csource{find_points.c}{1473}{1476}); the correction of the
  third stage by its cost rule from the counters, which starts at the new $s_2$
  (\csource{find_points.c}{1487}{1505}); and \verb|adapt_primes|: the words of the next correction
  set to twice the words swept so far, as a \verb|UInt64| that wraps, the rates, the second stage in
  the mode $2$ and more, and the third, every double in the order of the C code
  (\csource{find_points.c}{1383}{1388}, \csource{find_points.c}{1390}{1398},
  \csource{find_points.c}{1421}{1438}, \csource{find_points.c}{1477}{1486},
  \csource{find_points.c}{1506}{1509}, \csource{find_points.c}{1515}{1515}). The state of the
  correction: the counters of the run, the blocks folded, the decisions, the ends $s_2$ and $s_3$ of
  the last block folded and the latest decision with the words of the next correction, which the C
  code keeps in the arguments (\csource{find_points.c}{3755}{3762}); its start, nothing counted, the
  decision of the choice for block $0$ (\csource{find_points.c}{3853}{3856},
  \csource{find_points.c}{3861}{3863}); and the fold of the next block: its counters added, its ends
  $s_2$ and $s_3$ those of the last block folded, and, once the words reach those of the next
  correction, the correction, with a decision appended for the blocks from
  $j + \mathrm{adaptLag} + 1$ on when it changed an end (\csource{find_points.c}{4102}{4106},
  \csource{find_points.c}{4114}{4127}).
\end{definition}

The run (Definition~\ref{def:sieveRun}) becomes the loop over the blocks: from the start of the
correction, for each block in turn, one block (Definition~\ref{def:runBlock}) with the decisions so
far, the fold of its report, and a return after the block in which the callback stopped the search;
it returns the delivery with the ends $s_2$ and $s_3$ of the last block folded, which the search
writes back (\csource{find_points.c}{5120}{5120}, \csource{find_points.c}{5152}{5160},
\csource{find_points.c}{5163}{5163}).

\begin{lemma}[The correction keeps the decisions well formed]\label{lem:correctionWF}
  \uses{def:correction, def:decisionWF}
  \lean{Ratpoints.Sieve.AdaptUpInvariant, Ratpoints.Sieve.AdaptDownInvariant, Ratpoints.Sieve.AdaptThirdInvariant, Ratpoints.Sieve.adaptUp_spec, Ratpoints.Sieve.adaptDown_spec, Ratpoints.Sieve.adaptThird_spec, Ratpoints.Sieve.adaptSecond_spec, Ratpoints.Sieve.adaptPrimes_spec, Ratpoints.Sieve.RunCounts.sp2_fold, Ratpoints.Sieve.Correction.WF, Ratpoints.Sieve.Correction.wf_start, Ratpoints.Sieve.runFold_spec, Ratpoints.Sieve.runFold_decisions_spec, List.forall_mem_drop_of_forall_getElem!}
  \leanok
  For any values of the doubles: the walk up ends between its start and the number of the entries;
  the walk down ends at most at $s_2$, at least at $s_1$ when $s_1 \le s_2$, with every entry
  between counting no factor; the third stage ends between the new $s_2$ and the number of the
  entries; so the correction of a well-formed latest decision is well formed. A state of the
  correction is well formed when it has a decision, the first for block $0$, every decision well
  formed, and the latest decision and the end $s_2$ of the counters too; the start is well formed
  and the fold of a block whose decision is one of the state keeps it so. The fold adds one block
  folded and appends decisions only for the blocks from $\mathrm{folded} + \mathrm{adaptLag} + 1$
  on.
\end{lemma}

\begin{proof}
  \leanok
  The contracts of the loops state only what the natural numbers do, with the measures the entries
  left above and the position below; no proof reads a value of a double, also where the walks
  compare costs that may be NaN.
\end{proof}

\begin{lemma}[The correction in closed form]\label{lem:correctionAfter}
  \uses{lem:correctionWF, lem:runBlock, lem:counters}
  \lean{Ratpoints.Sieve.blockReport, Ratpoints.Sieve.correctionAfter, Ratpoints.Sieve.blockDecision, Ratpoints.Sieve.folded_correctionAfter, Ratpoints.Sieve.exists_toList_decisions_correctionAfter, Ratpoints.Sieve.size_pos_and_block_zero_correctionAfter, Ratpoints.Sieve.wf_correctionAfter, Ratpoints.Sieve.eq_blockReport}
  \leanok
  The report of block $i$ with the decisions so far, in closed form: the sum of the counts of the
  denominators of the block that pass the test of the shape (Lemma~\ref{lem:counters}), with the
  decision in force and its ends $s_2$ and $s_3$; the state of the correction after the folds of the
  first $k$ blocks with these reports, from the start; and the decision of block $i$, the decision
  in force for it after the folds of the blocks before it. They are functions of the arguments of
  the run that read neither the callback nor the flags of the output. After $k$ folds, $k$ blocks
  are folded, the decisions are those after $m \le k$ folds followed by decisions for blocks from
  $m + \mathrm{adaptLag} + 1$ on, and the state is well formed when the decision of the choice is; a
  block on which the callback does not stop the search reports what the closed form says.
\end{lemma}

\begin{proof}
  \leanok
  By induction on the number of folds, from Lemma~\ref{lem:correctionWF}; the report of a block from
  the contract of one block (Lemma~\ref{lem:runBlock}).
\end{proof}

\begin{theorem}[The lag lemma]\label{thm:lag}
  \uses{lem:correctionAfter, lem:runBlock}
  \lean{Ratpoints.Sieve.eq_of_lt_and_block_le, Ratpoints.Sieve.findDecision_eq_of_toList_eq_append, Ratpoints.Sieve.getElem!_findDecision_eq_of_toList_eq_append, Ratpoints.Sieve.getElem!_findDecision_correctionAfter}
  \leanok
  The decision in force for block $i$ among the decisions after the folds of the first $k$ blocks is
  the decision of block $i$ for every $k \ge i - \mathrm{adaptLag}$: block $i$ reads only the
  decisions made at the folds of the blocks before block $i - \mathrm{adaptLag}$, that is of the
  blocks at least $\mathrm{adaptLag} + 1$ before it.
\end{theorem}

\begin{proof}
  \leanok
  The decision in force is unique; decisions appended for blocks after $i$ change neither its index
  nor itself; and the decisions after $k$ folds are those after $i - \mathrm{adaptLag}$ folds
  followed by decisions for blocks from $i + 1$ on (Lemma~\ref{lem:correctionAfter}).
\end{proof}

\begin{theorem}[The block lemma]\label{thm:blocks}
  \uses{def:sieveRun, lem:correctionAfter, lem:runBlock, lem:blockDens, lem:deliver}
  Proposed names: \lname{Ratpoints.Sieve.blockPoints},
  \lname{Ratpoints.Sieve.SieveRunLoopInvariant}, \lname{Ratpoints.Sieve.sieveRunLoopInvariant\_nil},
  \lname{Ratpoints.Sieve.SieveRunLoopInvariant.pre},
  \lname{Ratpoints.Sieve.SieveRunLoopInvariant.quit},
  \lname{Ratpoints.Sieve.SieveRunLoopInvariant.append},
  \lname{Ratpoints.Sieve.SieveRunLoopInvariant.post}, \lname{Ratpoints.Sieve.sieveRun\_spec},
  \lname{Ratpoints.Sieve.fst\_run\_sieveRun},
  \lname{Ratpoints.Sieve.exists\_run\_sieveRun\_eq\_deliver}. For sieve entries whose factors have
  their tables of inverses and whose entries that count no factor are prime entries, the reciprocals
  of their moduli, the row offsets of the classes, the word of the classes that admit a numerator, a
  well-formed decision of the choice, forbidden divisors with their reciprocals, $A \ge 1$,
  $1 \le b_{\mathrm{low}}$, $b_{\mathrm{high}} < 2^{63}$ and $H < 2^{63}$, the run, with a callback
  that has not stopped the search and a store that fits all the entries, delivers the points of its
  blocks in order, each the points of the denominators of the block that pass the test of the shape
  with the first $s_3$ entries of the decision of the block, a function of the block and its
  decision, for every choice of the words of the classes and the ranges of the numerators. When the
  listing of the shape has no duplicates there is, for each denominator $b$, a number $s_3(b)$ of
  entries, the end of the third stage of the decision of its block, which depends neither on the
  callback nor on the flags of the output, such that the run delivers, over the denominators of the
  listing that pass the test of the shape, in order, the points of the check of the numerators in
  the ranges of $b$ that pass its class word and survive the first $s_3(b)$ entries.
\end{theorem}

\begin{proof}
  The invariant of the loop over the blocks: the delivery is that of the points of the blocks done;
  while the callback has not stopped the search, the state of the correction is that in closed form
  after the folds of the blocks done and the store fits all the entries. A stop in a block ends the
  run before any decision that the fold of that block pushes is read. The form per denominator
  follows from the block-listing lemma (Lemma~\ref{lem:blockDens}), each denominator lying in one
  block.
\end{proof}

\begin{theorem}[The blocks sifted apart]\label{thm:threads}
  \uses{thm:blocks, thm:lag, lem:runBlock, lem:deliver}
  Proposed names: \lname{Ratpoints.Sieve.deliverSlots},
  \lname{Ratpoints.Sieve.DeliverSlotsInvariant}, \lname{Ratpoints.Sieve.DeliverSlotsInvariant.quit},
  \lname{Ratpoints.Sieve.DeliverSlotsInvariant.append}, \lname{Ratpoints.Sieve.deliverSlots\_spec},
  \lname{Ratpoints.Sieve.toList\_info\_run\_runBlock\_appendPoint},
  \lname{Ratpoints.Sieve.run\_deliverSlots\_eq\_fst\_run\_sieveRun}. The property that the pool of
  threads of the C code preserves, which the model leaves out: block $i$ sifted alone into a slot
  from the empty list, with the decisions after the folds of any $k_i \ge i - \mathrm{adaptLag}$
  blocks and a store of its own that fits all the entries, gives exactly the points of block $i$;
  and delivering the slots of the blocks in order to the callback, stopping at its first stop, gives
  what the run delivers.
\end{theorem}

\begin{proof}
  From the lag lemma, the decision of each slot is that of its block, and the points of a block do
  not depend on the store that it starts from as long as it fits; the delivery of the slots in order
  is the delivery of the concatenation (Lemma~\ref{lem:deliver}).
\end{proof}

\section{The search with the choice}\label{sec:choiceSearch}

\begin{definition}[The search with the choice of the moduli]\label{def:chooseAndSearch}
  \uses{def:choice, def:sieving, def:pointsAtInfinity, def:delivery, def:report, def:sieveOptions}
  Proposed name: \lname{Ratpoints.chooseAndSearch}. The choice and the search after it, which the
  search of the polynomial (Definition~\ref{def:searchPoints}) calls after the 2-adic information:
  the choice of the moduli (Definition~\ref{def:choice}) and nothing delivered at its exit, after
  the message of the exit; the report of the choice and of its moduli; the points at infinity, after
  which the search returns when the callback has stopped it; the lines before the sieve, after which
  it returns when the Sturm step left no interval; and the sieving with the result of the choice
  (Definition~\ref{def:sieving}), returning the numbers of the moduli used, $s_1$ of the choice and
  the ends $s_2$ and $s_3$ of the last block, and $0$ for each after the returns before the sieve
  (\csource{find_points.c}{4996}{5006}, \csource{find_points.c}{5010}{5012},
  \csource{find_points.c}{5100}{5103}, \csource{find_points.c}{5105}{5106},
  \csource{find_points.c}{5210}{5213}).
\end{definition}

\begin{definition}[The search with the print channel]\label{def:findPointsVerbose}
  \uses{def:searchPoints, def:reversal, def:sturmRegion, def:Outcome, def:delivery,
    def:Options, def:sieveOptions, def:normalize, def:searchIntervals, def:useCLong, def:maxDegree}
  \lean{Ratpoints.findPointsWork, Ratpoints.findPointsVerbose}
  \leanok
  The search of \verb|find_points_work_1| after the test of the leading zeros, with the report
  (I.15): of the arguments, the reversal check, the search intervals, the Sturm step and the search,
  with whether the caller left the number of primes to the library, a refusal coming with the data
  of the callback as the report has left it (\csource{find_points.c}{4577}{4585}); and the search of
  \verb|find_points| with the callback and the print channel, the refusals of the program in its
  order and then that search, a refusal coming with the data of the callback
  (\csource{find_points.c}{4523}{4528}, \csource{find_points.c}{4546}{4546}, with I.13
  \verb|find_points.c:4551-4552|, \csource{find_points.c}{4557}{4559},
  \csource{find_points.c}{4564}{4565}, \csource{find_points.c}{4567}{4571},
  \csource{find_points.c}{4606}{4606}, \csource{find_points.c}{4624}{4633},
  \csource{find_points.c}{4658}{4658}, \csource{find_points.c}{5219}{5220},
  \csource{find_points.c}{5226}{5229}, \csource{find_points.c}{5235}{5235},
  \csource{find_points.c}{5241}{5242}, \csource{ratpoints.h}{362}{362},
  \csource{ratpoints.h}{399}{400}, \csource{ratpoints.h}{409}{411}, \csource{ratpoints.h}{427}{430},
  \csource{ratpoints.h}{436}{437}, \csource{ratpoints.h}{441}{442}).
\end{definition}

The interface keeps its shape. A text is printed through the print channel of a search, a function
that prints into the data of the callback, which holds the standard output of the caller, so that
the report comes between the points as on the standard output of the C program
(\lname{Ratpoints.Delivery.print}, which joins the delivery, Definition~\ref{def:delivery}, and
models no C code of its own). $\mathrm{findPointsWith}$ (Definition~\ref{def:findPointsWith}) takes
the options $s$ of the sieve for the parameters and a trailing print channel, by default the channel
that prints nothing, as without \verb|RATPOINTS_VERBOSE|, and is $\mathrm{findPointsVerbose}$ with
the data of the callback dropped from a refusal; it models no C code of its own, its links going to
$\mathrm{findPointsVerbose}$ and, the lines 4578-4585, to $\mathrm{findPointsWork}$.
$\mathrm{findPoints}$ (Definition~\ref{def:findPoints}) takes $s$, by default the options of a run
without options, and the appending callback, as before. The search of the polynomial
(Definition~\ref{def:searchPoints}) takes the normalized options, whether the caller left the number
of primes to the library and the print channel, and calls $\mathrm{chooseAndSearch}$ after the
2-adic information, its links of the choice and the sieve going there and its lines of the report to
the report (Definition~\ref{def:report}). The sieving (Definition~\ref{def:sieving}) takes the
options and the result of the choice and returns the ends $s_2$ and $s_3$ of the last block with the
delivery. The outcome (Definition~\ref{def:Outcome}) gains the numbers of the moduli used, $0$ by
default, its link \csource{ratpoints.h}{377}{383} in place of \csource{ratpoints.h}{377}{377}. The
certified statements use $\mathrm{findPointsWith}$ with the print channel at its default, so that
their text does not mention it.

\begin{lemma}[The search with the default print channel]\label{lem:choiceSearch}
  \uses{def:findPointsVerbose, def:findPointsWith}
  \lean{Ratpoints.findPointsWith_eq}
  \leanok
  With the print channel that prints nothing, the search with a callback is the refusals in the
  order of the program followed by the reversal check, the search intervals, the Sturm step and the
  search of the polynomial: the report prints nothing, and a refusal gives back its error alone.
\end{lemma}

\begin{proof}
  \leanok
  The print channel that prints nothing leaves a delivery as it is, by \verb|rfl|, and the error of
  a refusal is the first component of its pair.
\end{proof}

The lemmas of the search are restated on the choice (I.11). A sound run (Lemma~\ref{lem:soundRun})
no longer reads the end of the second stage: its sieve entries are sound, those that count no factor
are prime entries, with the reciprocals, the tables of inverses, the row offsets of the classes for
all the entries, the packings of the classes, the word of the classes and a store that fits all the
entries; a result of the choice that has $\mathrm{IsSievingInfo}$ gives one after the row offsets,
with a well-formed first decision (Lemma~\ref{lem:choice}, Definition~\ref{def:decisionWF}). The
contract of the run of the search (\verb|Ratpoints.exists_isSoundTest_run_sieveRun|,
Lemma~\ref{lem:searchList}) gives a sound test of the numerators block by block, $S(b, a)$ the test
of the sieve with the first $s_3(b)$ entries of the decision of the block of $b$
(Theorem~\ref{thm:blocks}), with the words that the packing gives, the ranges of the bounds, the
forbidden divisors of the choice and the listing of the search, which has no duplicates from
$1 \le b_{\mathrm{low}}$; its soundness is the survival of the points for any prefix of the entries
(Lemma~\ref{lem:survives}). The contract of the search
(\verb|Ratpoints.exists_isSoundTest_run_searchPoints|) is that of stage~H for normalized options,
$1 \le H < 2^{63}$, $1 \le b_{\mathrm{low}}$, $b_{\mathrm{high}} \le H$ and well-formed search
intervals, with the default print channel: there is a sound test, which depends neither on the
callback nor on the flags of the output, such that the search with any callback, flags and delivery
whose flag is clear gives the data and the count of the delivery of the points of the list of the
search (Definition~\ref{def:searchList}) for that test; an exit of the choice
(Lemma~\ref{lem:choice}) or the 2-adic exit gives a curve without points, for which the test that
lets nothing through is sound. The test depends on the search intervals, which the estimate of the
run reads and with it the choice and the decisions; the contract quantifies the test after them, as
in stage~H, and I.11 corrects the docstrings of both contracts, which still say otherwise. The
theorems downstream keep their statements but for the parameters: the reduction to the list of the
search under the exact check, where the survival of the entries drops out since the exact check
decides, the theorems of stage~D on the reversal, the order of stage~E, which goes through the
block-listing lemma (Lemma~\ref{lem:blockDens}), the protocol and the statements under
\verb|RATPOINTS_NO_CHECK| of stage~G and those of the search region of stage~H, now for every $s$
with no hypothesis on it.

\section{What comparator certifies}\label{sec:certifiedI}

\verb|comparator/Challenge.lean| keeps the specification that it restates and its imports; from the
model it no longer takes \verb|Params|, \verb|Params.WF| and \verb|sieveParams|. The eleven
statements, for all $c$, $H$, options $s$ of the sieve and options $o$, with $D$ the values and
$B_D$ the bits of the intervals of $o$, the search region that of $D$, $o.\mathrm{bLow}$ and
$o.\mathrm{bHigh}$, and $h_S$ the hypothesis $0 \le o.\mathrm{sturm} \to \mathrm{SturmTheorem}$:
\begin{enumerate}
  \item a triple lies in $\mathrm{bruteForce}(c, H)$ exactly when it is a normalized point of height
    at most $H$ (\verb|challenge_mem_bruteForce_iff|, unchanged);
  \item $\mathrm{bruteForce}(c, H)$ has no duplicates (\verb|challenge_nodup_bruteForce|,
    unchanged);
  \item under $h_S$, for $\mathrm{noCheck}$ false and $\mathrm{findPoints}(c, H, s, o) = L$:
    $p \in L$ exactly when $p$ is a normalized point of height at most $H$, with $y \ge 0$ if
    $\mathrm{noY}$ is set, in the search region
    (\verb|challenge_mem_iff_isPoint_of_findPoints_eq_ok|);
  \item for $\mathrm{findPoints}(c, H, s, o) = L$, $L$ has no duplicates
    (\verb|challenge_nodup_of_findPoints_eq_ok|);
  \item with the Sturm step on, $\mathrm{findPoints}(c, H, s, o)$ refuses exactly when $c$ has fewer
    than two or more than $\mathrm{maxDegree} + 1$ entries, $H = 0$, $H \ge 2^{63}$, $\Phi_c$ is not
    squarefree or $B_D$ is not well formed
    (\verb|challenge_exists_findPoints_eq_error_iff_of_sturm_nonneg|);
  \item with the Sturm step off, exactly when $c$ has fewer than two or more than
    $\mathrm{maxDegree} + 1$ entries, $H = 0$, $H \ge 2^{63}$, $d + 2 \le N$, or the reversal is
    allowed, the search region is the whole box, $d$ is even and $c_0 = c_1 = 0$, or $B_D$ is not
    well formed (\verb|challenge_exists_findPoints_eq_error_iff_of_sturm_neg|);
  \item under $h_S$, for $\mathrm{noCheck}$ false and $\mathrm{findPoints}(c, H, s, o) = L$: $L$ is
    the list of the points in the order of ratpoints for $\mathrm{maxDivisors}$, the switch of the
    reversal of $o$ and the whole box, $c$, $H$ and the largest denominator
    $\mathrm{denominatorBound}(o.\mathrm{bHigh}, H)$, filtered by $\neg\mathrm{noY} \lor y \ge 0$
    and the search region (\verb|challenge_eq_orderedPoints_of_findPoints_eq_ok|);
  \item for $\mathrm{noCheck}$ set and $\mathrm{findPoints}(c, H, s, o) = L$: every member of $L$ is
    $(x, 0, z)$ with $\gcd(x, z) = 1$, $z > 0$ or ($z = 0$ and $x = 1$), $|x| \le H$ and $z \le H$,
    in the search region (\verb|challenge_exists_eq_of_mem_of_findPoints_eq_ok_of_noCheck|);
  \item under $h_S$, for $\mathrm{noCheck}$ set and $\mathrm{findPoints}(c, H, s, o) = L$: for every
    normalized point $(x, y, z)$ of height at most $H$ in the search region, $(x, 0, z) \in L$
    (\verb|challenge_mem_of_isPoint_of_findPoints_eq_ok_of_noCheck|);
  \item for $\mathrm{noCheck}$ set, $\mathrm{findPoints}(c, H, s, o) = L$ and the search with
    $\mathrm{noCheck}$ cleared and $\mathrm{noY}$ set returning $L'$: $L'$ with $y$ replaced by $0$
    is a sublist of $L$ (\verb|challenge_map_sublist_of_findPoints_eq_ok_of_noCheck|);
  \item for every callback with its data, $\mathrm{findPointsWith}(c, H, s, o)$ with them gives,
    mapped to its data and count, what the delivery of the list of $\mathrm{findPoints}(c, H, s, o)$
    to the callback gives, and the same refusal, with the same error, exactly where
    $\mathrm{findPoints}$ refuses (\verb|challenge_map_findPointsWith_eq_map_findPoints|).
\end{enumerate}
Statements 3 to 11 are those of stage~H with $s$ in place of $P$ and without $P.\mathrm{WF}$;
\verb|challenge_wf_sieveParams| goes. Each is stronger than before: it holds for the program's own
choice of the moduli and its corrections, for every value of the options of the sieve, the tuning
values included, and whatever values the doubles of the cost model take. Statement~10 compares two
searches with the same options $s$, which choose the same moduli and make the same decisions, since
the decisions depend neither on the callback nor on the flags of the output, so that the survivors
under $\mathrm{noCheck}$ are those of the same test. Statement~11 holds for the default print
channel; the report under \verb|-v| is not proved (Section~\ref{sec:moduliReport}). The challenge's
module docstring and \verb|comparator/README.md| record the change by stage, with the earlier
acceptances as history (I.10).

\section{The numbers of the moduli used and the report under -v}\label{sec:moduliReport}

The numbers of the moduli used are a field of the outcome (Definition~\ref{def:Outcome}), written as
Section~\ref{sec:moduliC} says; the difference $s_3 - s_2$ that message~4 prints is an integer, for
the \verb|long|. The report of the program (Definition~\ref{def:cliReport}) prints message~4, its
link becoming \csource{main.c}{979}{1033} in place of the two links around the lines of message~4.

\begin{definition}[The report of the library]\label{def:report}
  \uses{def:Double, def:reversal, def:setupUs1, def:numClasses, def:choice, def:forbidden}
  Proposed names: \lname{Ratpoints.Report.width3}, \lname{Ratpoints.Report.longs},
  \lname{Ratpoints.Report.arguments}, \lname{Ratpoints.Report.reversal},
  \lname{Ratpoints.Report.sturmStep}, \lname{Ratpoints.Report.setup},
  \lname{Ratpoints.Report.noClasses}, \lname{Ratpoints.Report.packing},
  \lname{Ratpoints.Report.noPoints}, \lname{Ratpoints.Report.choice},
  \lname{Ratpoints.Report.excludedValuations}, \lname{Ratpoints.Report.excluded},
  \lname{Ratpoints.Report.moduli}, \lname{Ratpoints.Report.atInfinity},
  \lname{Ratpoints.Report.sieving}. The texts of the report that the library prints under
  \verb|RATPOINTS_VERBOSE|, each printed through the print channel where the C code prints it, or
  right after a pure computation when nothing else prints in between: the conversions \verb|%3ld|
  and \verb| %ld| in loops; the arguments as the search normalizes them
  (\csource{find_points.c}{4691}{4725}); the outcome of the reversal check, its case read from the
  degree and $c_0$ as the check decides (\csource{find_points.c}{4740}{4743},
  \csource{find_points.c}{4757}{4758}, \csource{find_points.c}{4771}{4774},
  \csource{find_points.c}{4791}{4801}, \csource{find_points.c}{4804}{4807}); the Sturm step
  (\csource{find_points.c}{4860}{4861}, \csource{find_points.c}{4864}{4878}); the set-up of the
  denominators (\csource{find_points.c}{4897}{4898}, \csource{find_points.c}{4915}{4918},
  \csource{find_points.c}{4921}{4926}, \csource{find_points.c}{4928}{4941},
  \csource{find_points.c}{4946}{4947}); the 2-adic exit (\csource{find_points.c}{4960}{4961}); the
  packing of the classes and the number of primes to look at (\csource{find_points.c}{4969}{4985},
  \csource{find_points.c}{4988}{4991}); the exit of the choice
  (\csource{find_points.c}{5007}{5009}); the bits per word, the primes looked at and the cut points
  of the choice (\csource{find_points.c}{3089}{3095}); the forbidden valuations and divisors
  (\csource{find_points.c}{5049}{5058}, \csource{find_points.c}{5042}{5048},
  \csource{find_points.c}{5059}{5063}); the moduli of the three stages and the cost of a check,
  \verb|%.0f| and \verb|%.2f| of \verb|Float|s (\csource{find_points.c}{5027}{5041},
  \csource{find_points.c}{5064}{5066}); the points at infinity
  (\csource{find_points.c}{5075}{5077}); and the sieve (\csource{find_points.c}{5104}{5104},
  \csource{find_points.c}{5111}{5113}). The doubles are printed as \verb|printf| of glibc prints
  them under \verb|%.pf|, by \lname{Ratpoints.Double.printFixed}, which joins
  Definition~\ref{def:Double}, since the \verb|%f| of the program prints with it too: the exact
  value rounded to $p$ places, ties to even, in full, without the point for $p = 0$, after a
  \verb|-| whenever the sign bit is set.
\end{definition}

The ranges that these texts take were excluded as instrumentation, or cut out of the components that
run around them: the reversal check (Definition~\ref{def:reversal}) keeps
\csource{find_points.c}{4732}{4739}, \csource{find_points.c}{4744}{4756},
\csource{find_points.c}{4759}{4770}, \csource{find_points.c}{4775}{4790} and
\csource{find_points.c}{4802}{4803}; the Sturm step of the search (Definition~\ref{def:sturmRegion})
\csource{find_points.c}{4857}{4859}, \csource{find_points.c}{4862}{4863} and
\csource{find_points.c}{4879}{4892}; the points at infinity (Definition~\ref{def:pointsAtInfinity})
\csource{find_points.c}{4894}{4896}, \csource{find_points.c}{4899}{4899},
\csource{find_points.c}{5068}{5074} and \csource{find_points.c}{5078}{5099}. The parts of the
excluded ranges under \verb|DEBUG| stay excluded: 4727-4730, 4809-4812, 4965-4968, 5014-5026 and
5108-5110 of \verb|find_points.c|. The report is tested byte for byte, not proved: with the print
channel that prints nothing every text vanishes (Lemma~\ref{lem:choiceSearch}), and the theorems
speak about that channel.

\section{The program}\label{sec:moduliCli}

\begin{definition}[The tuning options and the report in the program]\label{def:cliTuning}
  \uses{def:sieveOptions, def:cliNumbers, def:cliEnds, def:cliProcess}
  Proposed names: \lname{Ratpoints.SieveOptions.store}, \lname{Ratpoints.Cli.Data.print}. The
  options of the sieve with the field that an option of the command line sets: the C \verb|long| of
  its argument for \verb|-p -n -N -R -A -P -F| and its double for \verb|-r -U -C -Q -W|, as
  \verb|main.c| reads them, the Jacobi test off for \verb|-j| and on for \verb|-J|; it models no C
  code of its own, being the part of the reading of the options (Definition~\ref{def:cliInput}) that
  sets the fields of the options of the sieve. And the print channel of the program under \verb|-v|:
  the text appended to the output held in the data of the callback, for the calls of \verb|printf|
  in the library.
\end{definition}

The program (Definition~\ref{def:cliRun}) runs $\mathrm{findPointsVerbose}$ with its options of the
sieve, the array size $\mathrm{defaultArraySize}$ that it sets before the search
(\csource{main.c}{221}{222}, from the parameters of the executable with I.10), and the print channel
of \verb|-v| when the setting is on, which \verb|-q| clears; on a refusal it prints the head of the
report, the text that the report of the library has printed and the message of the error. The
options \verb|-r -R -U -C -A -Q -W| no longer go unread: they set the options of the sieve, with the
validation of \verb|main.c| as before; only \verb|-t| is checked and ignored, for good.

\begin{lemma}[The program without -v]\label{lem:cliVerbose}
  \uses{def:cliTuning, def:findPointsWith, def:cliInput, def:cliFailure}
  Proposed names: \lname{Ratpoints.Cli.findPointsVerbose\_eq\_ok\_of\_findPointsWith\_eq\_ok},
  \lname{Ratpoints.Cli.verbose\_eq\_false\_of\_input\_eq\_ok}. A search with the print channel that
  prints nothing that gives an outcome gives the same outcome as $\mathrm{findPointsVerbose}$ with
  that channel; and a command line that the program reads with \verb|-q| has the setting of
  \verb|-v| clear.
\end{lemma}

\begin{proof}
  \uses{thm:cliRegion}
  By the definition of $\mathrm{findPointsWith}$ on the two cases of the result, and by the last
  conjunct that the reading of a command line gains, that \verb|-q| cancels \verb|-v|.
\end{proof}

The theorem on the program (Theorem~\ref{thm:cliRun}) gains the hypothesis that the setting of
\verb|-v| is clear, and speaks about $\mathrm{findPointsWith}$ and $\mathrm{findPoints}$ with the
options of the sieve of the command line and the program's array size, with no hypothesis of
well-formedness; its corollaries under \verb|-q| get the hypothesis from Lemma~\ref{lem:cliVerbose}.
Message~4 comes in the report over the outcome, as messages~7 and~8 do. The reading of a command
line (Theorem~\ref{thm:cliRegion}) gains that \verb|-q| cancels \verb|-v|, and its proof covers the
branches of the tuning options. The report under \verb|-v| is not proved
(Section~\ref{sec:moduliReport}), and the program theorem is not certified, as before.

\section{The differential tests}\label{sec:moduliTests}

The harness compares the choice, the blocks and the report with the C program, which must be the
default build at the pin, without a \verb|tuning.mk|, the file of machine-dependent tuning constants
that \verb|make tune| writes and the Makefile includes (the harness refuses a build directory that
has one), and without flags that tune the choice of the moduli. The comparisons that depend on the
choice, the pairs under \verb|-x|, message~4 and the report under \verb|-v|, hold on one machine
(Section~\ref{sec:costs}); CI compares them on its runner. The items of
Section~\ref{sec:moduliFiles} bring these comparisons in stages, each on the runs whose output it
makes those of ratpoints; in the end:
\begin{itemize}
  \item under \verb|-x| the pairs are compared exactly with the C program on the same command line,
    on every run, in place of the comparison with the moduli pinned and the weaker check alone where
    the C program chose other moduli; the weaker check of the survivors stays as an extra check,
    except under \verb|-v|, where the report stands between the pairs; the generators append
    \verb|-p 30 -N 30| to a run under \verb|-x| at times;
  \item message~4 is compared in every report; the runs of \verb|test4.sh| whose record drops it are
    compared with the live C program, a disagreement of the C program with its record reported
    apart;
  \item the report under \verb|-v| is compared byte for byte, the runs \verb|v()| of \verb|test4.sh|
    with the live C program, its filters of the bits per word and of the iterations and the piped
    lines of \verb|test3.sh| on the whole output before the pipe;
  \item a new default suite \verb|tuning| runs $300$ random values of \verb|-r -R -U -C -A -Q -W|,
    among them $\infty$, $0$, the least subnormal and huge values, at times with the cut points and
    the numbers of primes given, at heights up to $3000$;
  \item a new default suite \verb|threads| runs the C program with \verb|-t 4| and \verb|-t 7|
    against \verb|ratpoints-lean|, and a new default suite \verb|verbose| $300$ random runs under
    \verb|-v|; the errors suite gains the cases of the report: each exit, error~8 after the report,
    the three cases of the reversal, the divisors and the failed factorization, the valuations
    excluded and the tuning options;
  \item the high suite gains runs at $16383$ where the correction takes effect, under \verb|-x|
    without \verb|-A|, with \verb|-A 1| and with \verb|-A 2|, and one with \verb|-z| (message~4); a
    run at $16383$ with \verb|-t 4| under \verb|-x| joins the suite \verb|threads|;
  \item a local suite \verb|adapt| runs $23$ command lines without \verb|-A| and with \verb|-A 0|,
    \verb|-A 1| and \verb|-A 2| under \verb|-x|, also above the largest height of the default
    suites, where the correction may take effect; and a local pair compares a build of the C program
    with \verb|-DRP_BLOCK_LEN=3| against the model with $\mathrm{blockLen} = 3$, on the suites that
    do not compare with recorded outputs.
\end{itemize}
In the prototype the default suites compare $5397$ times with none failed: \verb|test4| $590$,
\verb|test3| $21$, \verb|rptest| $1206$, \verb|random| $300$, \verb|outputs| $300$, the errors suite
$384$, the fuzz $800$, the parameters $272$, the high suite $19$, \verb|regions| $400$,
\verb|depths| $300$, \verb|tuning| $300$, \verb|threads| $205$ and \verb|verbose| $300$; $462$ runs
under \verb|-x| are compared exactly ($17$ of them where the C program may correct its moduli),
$445$ reports carry message~4 and $414$ runs are under \verb|-v| with the report of the library.
\verb|test4high| compares $13$, the suite \verb|adapt| $92$ under \verb|-x| ($69$ where a correction
is possible; in the C program $18$ of the $23$ command lines change their pairs with \verb|-A|), and
the pair with the block length $3$ compares $3601$ times in twelve default suites and $92$ times in
\verb|adapt|, none failed. Every point list is compared in its order.

The \verb|#guard| examples check the bits of the \verb|Float| constants against gcc, \verb|%.pf|
against glibc, and the parts of the choice that do not depend on the platform's \verb|log| and
\verb|pow| on small pools: the candidates that the stages take and drop, the composite moduli
offered and refused (an example is $15$ refused while $45 = 9 \cdot 5$ is offered, since the prime
power $9$ says something), the ties of the least density, the exit of the third stage and the stable
sort of the end of the choice. No guard pins an output that depends on the choice of the moduli:
under \verb|-x| the guards check that the pairs contain those of the points and are normalized pairs
in the region, message~4 is masked, and of the report under \verb|-v| only the text before the
moduli and after the cost of a check is compared.

\section{Deviations and exclusions}\label{sec:moduliDeviations}

Each component differs from the C code in a stated way, recorded in its \verb|Deviation:| paragraph:
\begin{itemize}
  \item the doubles of the cost model as Lean's \verb|Float| with the arithmetic of the platform,
    \verb|log| and \verb|pow| the functions of the platform's C library that the Lean runtime binds,
    the same routines as ratpoints' on every argument that occurs (Section~\ref{sec:costs});
    $\log 2$ computed at run time, where gcc folds \verb|log(2.0)|; \verb|ceil| and \verb|floor|
    calls of the library where gcc compiles exact instructions, and the number of set bits of a word
    a loop, both recorded from I.3 in the \verb|Deviation:| paragraphs of the components that use
    them; the tuning values and the ends of the intervals entering by their bits;
  \item the sorts by \verb|qsort| as Lean's stable merge sort with the test whether the comparison
    of the C code is at most $0$, which gives the order of glibc's \verb|qsort| on keys that are
    finite or $+\infty$;
  \item the examination of a prime returning the code and the density with its table, the exit as
    none and the density $0.0$ for the code $0$, where the C code leaves it unset; the candidates as
    records returned, their key and cost $0.0$ until the caller sets them; the entry of a prime
    power made for each composite modulus that needs it, where \verb|make_modulus| keeps it by its
    index; no test of room for the prime powers, which never fails; the loop of the tables a loop
    over the positions that ends at the number to look at, which the body raises;
  \item the table of rows of a prime of the third stage made after the choice, and none for a prime
    that the third stage examines and passes over (lines 2994-3000), where \verb|examine_prime|
    makes it with the entry;
  \item the counters a record of their own passed beside the delivery and the store, for the six
    fields of the sieving thread, \verb|UInt64| for the \verb|unsigned long|s; the report of a block
    returned with the delivery and the store, for \verb|worker_stat|; the counters zeroed by
    \verb|worker_init| (\csource{find_points.c}{3127}{3128}) part of the block, which zeroes them
    again;
  \item the decisions one array of records for three arrays; the fields $s_2$, $s_3$ and the words
    of the next correction kept in the state of the correction, where \verb|adapt_primes| writes
    them into the arguments, since only the run reads them; \verb|RP_PUBLISH| and \verb|RP_ACQUIRE|
    a plain store and a plain read; the entries and the reciprocals of each block cut as lists;
  \item \verb|RP_BLOCK_LEN| an optional constant of the model; the shapes' loops taking the
    denominator's step as an argument, for the two calls of each shape;
  \item \verb|RATPOINTS_VERBOSE| as the print channel of the search, which prints into the data of
    the callback, a text printed where the C code prints it or right after a pure computation; a
    refusal of $\mathrm{findPointsVerbose}$ returning the data of the callback with its error;
  \item the outcome writing back the numbers of the moduli used, the flags and the search intervals
    and nothing else: not \verb|args->sp3|, which \verb|find_points_work| zeroes and does not
    restore and the last block sets (lines 4551, 4561 and 5163), nor the private fields
    \verb|adapt_at|, \verb|run_words|, \verb|run_denoms|, \verb|sp3_max| and \verb|check_rel|;
  \item the difference $s_3 - s_2$ of message~4 an integer.
\end{itemize}
Left out, with their reasons in the record of the correspondence: the threads (kind threads), the
pool, its workers, the collection of the points of a block and \verb|run_threaded| with their stubs
(\csource{find_points.c}{4129}{4519}), the declarations of the lock of the tables and of the
destruction of the pool (\csource{find_points.c}{41}{43}), the atomic loads and stores
(\csource{rp-private.h}{76}{84}), the pool, slot and stop of a thread
(\csource{rp-private.h}{584}{584}) and the fields of the number of threads and of the pool of the
arguments (\csource{ratpoints.h}{370}{376}, \csource{ratpoints.h}{398}{398}), since the pool gives
the output of the sequential loop (Theorem~\ref{thm:threads}); \verb|RP_PRIME_STATS| in
\verb|adapt_primes|, a probe that prints each correction to the standard error
(\csource{find_points.c}{1510}{1514}, kind instrumentation); \verb|RATPOINTS_DEFAULT_SP1|
(\csource{ratpoints.h}{42}{42}, kind other), read nowhere; and the private fields of the arguments
that hold memory, the scratch integers and the buffers of the entries, rows, integers, forbidden
words and prime powers (\csource{ratpoints.h}{385}{388}, \csource{ratpoints.h}{392}{394}, kind
memory). The private fields that hold data of the model go to their owners: the sieve entries, the
reciprocals and the forbidden divisors to the result of the choice, together with
\verb|stage3_list|, the buffer of the table of the third stage, which the model replaces by the
thunk of the checks (\csource{ratpoints.h}{389}{389}, \csource{ratpoints.h}{391}{391}); the
valuation records and the divisors to the divisor set-up (\csource{ratpoints.h}{390}{390},
Definition~\ref{def:setupUs1}); the estimates and the reserve to $\mathrm{Params}$
(\csource{ratpoints.h}{395}{397}); the end of the arguments to the search
(\csource{ratpoints.h}{399}{400}); the numbers of the moduli used to the outcome
(\csource{ratpoints.h}{377}{383}); the counters to their record (\csource{rp-private.h}{579}{581});
the decision of a block to the block (\csource{rp-private.h}{582}{582}); and the type of the sieving
thread to one denominator (\csource{rp-private.h}{585}{585}). With them every line of code of the C
files at the pin is recorded.

\section{What changes in the earlier chapters}\label{sec:moduliEarlier}

Each change is made by the item of Section~\ref{sec:moduliFiles} named in parentheses, the item that
makes the old text untrue, so that the blueprint stays true between the commits; a node that loses
names keeps its label, a node enters with the uses of the declarations that it holds at its item and
gains the others with the items that bring them, and the uses and the references of a deleted node
go to the node that takes its place, in the item that deletes it. The references to the labels that
I.12 deletes, Definitions~\ref{def:sieveParams}, \ref{def:examinePrimes}, \ref{def:modulusFactors},
\ref{def:makeModulus} and~\ref{def:sievingInfo} and Lemmas~\ref{lem:examinePrimes},
\ref{lem:modulusFactors} and~\ref{lem:sievingInfo}, are listed below with the item that retargets or
drops each; by the lines of the chapters as they stand before stage~I, they are: in
\verb|01-spec.tex| lines 300, 382, 403 and 428; in \verb|02-sieve.tex| lines 87, 97 and 446; in
\verb|03-words.tex| lines 137 and 496; in \verb|04-setup.tex| lines 452, 472 and 503; in
\verb|05-denominators.tex| lines 378 and 649; in \verb|06-tables.tex| lines 35, 395, 429, 435, 438,
464, 465, 471, 565, 588 and 600, besides the deleted nodes themselves; in \verb|07-stages.tex| lines
32, 113, 465 and 474; in \verb|08-packing.tex| lines 524 and 644; in \verb|09-rows.tex| lines 695,
697, 699, 704, 713 and 893; in \verb|10-output.tex| lines 577, 1102 and 1139; in
\verb|11-region.tex| line 863; and in this chapter the paragraph after Definition~\ref{def:choice},
the proof of Lemma~\ref{lem:choice} and this section.
\begin{itemize}
  \item The introduction of the blueprint: the proofs show that the model meets the specification
    for every value of the options of the sieve, whatever moduli the search chooses, no longer for
    every well-formed choice of the sieving parameters (I.10); the item of stage~I in the list of
    the stages names the choice of the moduli and the cut points by the curve, the blocks with the
    correction between them, the numbers of the moduli used, the report under \verb|-v| and the
    tuning options, the threads left out with the property that they preserve proved
    (Chapter~\ref{chap:moduli}) (I.16).
  \item Chapter~\ref{chap:spec}: the section on what comparator certifies
    (Section~\ref{sec:comparator-ab}) says that since stage~I the statements take the options of the
    sieve with no hypothesis on them, that \verb|wf_sieveParams| is gone and eleven statements
    remain, and that \verb|Params|, \verb|Params.WF| and \verb|sieveParams| are no longer taken from
    the model (I.10), with the record of each acceptance in the item that runs comparator (I.10,
    I.11, I.13 and I.15); the list of Section~\ref{sec:cli}: the program prints message~4 (I.13) and
    under \verb|-v| the report of the library between the points (I.15), and its item on the tuning
    options refers to the options of the sieve (Definition~\ref{def:sieveOptions}) for
    Definition~\ref{def:sieveParams} (I.12) and reads \verb|-r -R -U -C -A -Q -W| into the options
    of the sieve, only \verb|-t| being checked and ignored (I.14); the command line
    (Definition~\ref{def:cliInput}) refers to Definition~\ref{def:sieveOptions} for
    Definition~\ref{def:sieveParams}, its uses likewise (I.12), and sets the options of the sieve of
    the tuning options through \lname{Ratpoints.SieveOptions.store}, its uses gaining
    Definition~\ref{def:cliTuning} (I.14); the program (Definition~\ref{def:cliRun}) runs
    $\mathrm{findPointsVerbose}$ with the options of the sieve and the array size that it sets, its
    link \csource{main.c}{221}{222} added and its uses gaining
    Definition~\ref{def:findPointsVerbose} for Definitions~\ref{def:findPointsWith}
    and~\ref{def:sieveParams} (I.10), with the print channel of \verb|-v| (I.15); the section on the
    differential tests: its list of what the later stages added gains stage~I, the suite
    \verb|adapt| and the pair with the block length $3$ (I.11), the suites \verb|tuning| and
    \verb|threads| (I.14) and \verb|verbose| (I.15); the runs left out, the four of \verb|test4.sh|
    whose filter looks at the report of the library and those of \verb|test3.sh| with a filtered
    output, compared with message~4 (I.13) and the report (I.15); its closing paragraph: the pairs
    under \verb|-x| compared exactly (I.11), message~4 no longer dropped (I.13), the output under
    \verb|-v| byte for byte (I.15); and the counts, with the items that change them.
  \item Chapter~\ref{chap:sieve}: the introduction says that the theorems hold for every choice of
    the moduli that the search makes (I.16); the primes of the sieve
    (Definition~\ref{def:primeTable}) gain \lname{Ratpoints.getElem!\_primeTable\_le\_of\_le}
    (I.12); the parameters (Definition~\ref{def:Params}) gain $s_3$, the primes looked at, the
    estimates of the run, the cost of a check and the words of the first correction, their link
    \csource{ratpoints.h}{395}{397} added (I.8), are internal to the search (I.10) and the result of
    the choice (I.11), and lose $\mathrm{Params.WF}$, its instance and its three lemmas, with the
    references to Definitions~\ref{def:examinePrimes} and~\ref{def:modulusFactors} and
    Lemma~\ref{lem:modulusFactors}, the loop of the tables becoming Definition~\ref{def:tablesLoop},
    and their uses of Definitions~\ref{def:barrett} and~\ref{def:primeTable}, which only the
    well-formedness read (I.12); the parameters of the executable (Definition~\ref{def:sieveParams})
    give $\mathrm{sp2Extra}$, $\mathrm{defaultNumPrimes}$ and $\mathrm{defaultMaxForbidden}$ with
    their links \csource{ratpoints.h}{95}{97}, \csource{ratpoints.h}{333}{340} and
    \csource{ratpoints.h}{345}{352} to Definition~\ref{def:choiceConstants}, their uses gaining it
    (I.1), give their lines of the choice to its components (I.3, I.7, I.8, as the paragraph after
    Definition~\ref{def:choice} lists them), give $\mathrm{SieveOptions}$ with its links
    \csource{main.c}{337}{338}, \csource{main.c}{344}{344}, \csource{main.c}{347}{347},
    \csource{main.c}{349}{350} and \csource{main.c}{352}{352} to Definition~\ref{def:sieveOptions},
    where they become \csource{main.c}{337}{347}, \csource{main.c}{349}{350} and
    \csource{main.c}{352}{352}, their uses gaining it (I.8), give
    \csource{find_points.c}{4577}{4577} and \csource{main.c}{221}{222} to the search and the program
    (I.10), and go (I.12): $\mathrm{sieveParams}$, $\mathrm{defaultParams}$, $\mathrm{defaultSp1}$,
    \verb|wf_sieveParams| and \verb|wf_defaultParams| are deleted, with the exclusion of
    \csource{ratpoints.h}{42}{42}; the paragraph after it gives way to one on the choice made inside
    the search, as in the C code, and the theorems for every choice it makes: its sentences on the
    statements, on the quantification over $\mathrm{Params.WF}$ and on \verb|wf_sieveParams| (I.10),
    the rest, the model taking the choice as an input and the stages of ratpoints and of the
    parameters of the executable (I.11); the sentence of Section~\ref{sec:sieve} on the components
    that may run \verb|Id.run| names $\mathrm{findPointsVerbose}$ too (I.10); the run
    (Definition~\ref{def:sieveRun}) gives the links \csource{find_points.c}{3900}{3904} to the plan
    of the blocks and \csource{find_points.c}{3917}{3918}, \csource{find_points.c}{3926}{3936},
    \csource{find_points.c}{3954}{3966}, \csource{find_points.c}{4009}{4009},
    \csource{find_points.c}{4042}{4052}, \csource{find_points.c}{4059}{4059},
    \csource{find_points.c}{4066}{4072} and lines 4075-4076 to the loops of the shapes (I.5), lines
    3889-3898, 3911, 3913 and 4077-4078 to one block, line 3899 to no declaration (I.6), and line
    3861 to the correction (I.9); it is the loop over the blocks with the correction, its
    \verb|\lean| list $\mathrm{sieveRun}$ alone, its contract going to Theorem~\ref{thm:blocks}, its
    links \csource{find_points.c}{5120}{5120}, \csource{find_points.c}{5152}{5160} and
    \csource{find_points.c}{5163}{5163}, and its uses Definitions~\ref{def:runBlock}
    and~\ref{def:correction} (I.11); the points at infinity (Definition~\ref{def:pointsAtInfinity})
    give their lines of the report to Definition~\ref{def:report}, their links becoming
    \csource{find_points.c}{4894}{4896}, \csource{find_points.c}{4899}{4899},
    \csource{find_points.c}{5068}{5074} and \csource{find_points.c}{5078}{5099} (I.15); the search
    for the points (Definition~\ref{def:findPoints}) takes the options of the sieve (I.10) and, in
    place of the tables and the sieve entries of Definition~\ref{def:sievingInfo}, the choice of the
    moduli (Definition~\ref{def:choice}) and the run over the blocks (I.11); the list of the
    deviations after it: the listing of the denominators and the cut points fixed for the run
    (I.11), no output under \verb|VERBOSE| (I.15); the reduction of the run to the reference
    (Theorem~\ref{thm:run_reference}) loses \verb|Ratpoints.Sieve.run_sieveRun_of_eq_filter| and
    keeps \verb|List.flatMap_filter_flatMap_filter|, by which the contract of the run of the search
    moves the test of a filtered list into the test of the numerators: its statement becomes that
    fact on lists and its uses go (I.11); the theorem on the points (Theorem~\ref{thm:findPoints})
    takes the options of the sieve, its uses losing Definition~\ref{def:Params} and gaining
    Definition~\ref{def:sieveOptions} (I.10), its proof going through Lemma~\ref{lem:choice} (I.11);
    the paragraph after it: the statements take $s$, the eleven of them, \verb|wf_sieveParams| gone,
    the names that the challenge takes from the model and the record of the acceptances (I.10, the
    records with the items that run comparator), the contract of the run named in the file of the
    run (I.11); and the files section: the well-formedness in \verb|Ratpoints/Proof/Params.lean|
    (I.12), the run in \verb|Ratpoints/Model/Blocks.lean| and its contract in
    \verb|Ratpoints/Proof/Correction.lean| (I.11).
  \item Chapter~\ref{chap:words}: the array size (Definition~\ref{def:arraySize}) is a field of the
    options of the sieve, normalized by \lname{Ratpoints.arraySizeAt} with
    \lname{Ratpoints.one\_le\_arraySizeAt}, which take the links of $\mathrm{Params.arraySizeAt}$
    (I.8); the passage after it, on $\mathrm{sieveParams}$, the quantification over the field and
    the challenge, gives way to one on the field of the options of the sieve (I.10); the walk
    (Definition~\ref{def:walkWord}), the scan of a piece (Definition~\ref{def:scanPiece}), one piece
    (Definition~\ref{def:siftPiece}) and one denominator (Definition~\ref{def:sift}) take and return
    the counters, one denominator with the link \csource{rp-private.h}{585}{585} and the third
    conjunct of its contract, Lemma~\ref{lem:sift_eq_checkDen} takes the counters, the uses of the
    four definitions and of the lemma gaining Definition~\ref{def:counters}, and those of one
    denominator and of the proof of the lemma Lemma~\ref{lem:counters} too, and the deviation that
    leaves the counts to stage~I is reworded (I.4); the passage after the run on what the search
    hands it, the entries of Definition~\ref{def:sievingInfo}, is on the choice and the blocks
    (I.11); and the files section: \verb|Ratpoints/Proof/Sift.lean| split (I.4), the run and its
    contract (I.11).
  \item Chapter~\ref{chap:setup}: the examination of a prime with the exit
    (Definition~\ref{def:examinePrimeExit}) returns the code and the density, its link
    \csource{find_points.c}{1761}{1781} in place of \csource{find_points.c}{1761}{1765} (I.2); the
    reversal check (Definition~\ref{def:reversal}) gives its lines of the report to
    Definition~\ref{def:report}, its link becoming \csource{find_points.c}{4732}{4739},
    \csource{find_points.c}{4744}{4756}, \csource{find_points.c}{4759}{4770},
    \csource{find_points.c}{4775}{4790} and \csource{find_points.c}{4802}{4803} (I.15); the search
    of the polynomial (Definition~\ref{def:searchPoints}) gives up its link
    \csource{find_points.c}{2793}{2795}, lines 2794-2795 going to Definition~\ref{def:choice} and
    line 2793 to no declaration (I.8), takes the normalized options of the sieve, whether the caller
    left the number of primes to the library and the print channel, its uses following (I.10), calls
    the search with the choice, its links of the choice and the sieve (lines 4996-5006, 5010-5012,
    5100-5103 and 5105-5106) going to Definition~\ref{def:chooseAndSearch}, its text on
    Definition~\ref{def:sievingInfo} and Lemma~\ref{lem:sievingInfo} rewritten for the choice and
    Lemma~\ref{lem:choice} and its uses Definition~\ref{def:chooseAndSearch} for
    Definitions~\ref{def:Params}, \ref{def:sievingInfo} and~\ref{def:sieving} (I.11), and gives its
    lines of the report (4915-4918, 4921-4926, 4928-4941, 4960-4961, 5007-5009 and 5104) to
    Definition~\ref{def:report} (I.15); the paragraph after it on
    $\mathrm{findPointsWith}(c, H, P, o)$ takes $s$ (I.10); Theorems~\ref{thm:findPoints_perm}
    and~\ref{thm:refusals} take $s$, the proof of the refusal theorems using
    Lemma~\ref{lem:choiceSearch} (I.10); the deviation that leaves the codes of \verb|examine_prime|
    to the choice goes (I.2), the one on the arguments of \verb|sieving_info| that serve the choice
    goes (I.11), and the one on \verb|RATPOINTS_VERBOSE| is reworded (I.15); the deviation that
    names \verb|makeModulus?| and the closing passage on what belongs to stage~I are reworded
    (I.16); and the files section: the flagged contracts in \verb|Ratpoints/Proof/Sift.lean| (I.4),
    that of the run (I.11).
  \item Chapter~\ref{chap:denominators}: the introduction, which leaves the blocks to stage~I, says
    that the model has them (I.11); the forbidden divisors (Definition~\ref{def:forbidden}) keep the
    reciprocal of their prime, their uses gaining Definition~\ref{def:mulMod} (I.2), their reference
    to Definition~\ref{def:examinePrimes} going to Definition~\ref{def:tablesLoop} (I.12); the
    divisor set-up (Definition~\ref{def:setupUs1}) gains the link \csource{ratpoints.h}{390}{390}
    (I.8); the shapes (Definition~\ref{def:Shape}) give the lines 3750-3754 to the blocks and the
    lines of the listing (3812-3837, 3919-3925, 3937-3943, 3967-3991, 3993-4008, 4012, 4053-4058,
    4060-4065 and 4073-4074) to the plan of the blocks and the loops of the shapes (I.5), and the
    lines 3755-3762 to the correction (I.9), their link \csource{find_points.c}{3736}{3763} becoming
    \csource{find_points.c}{3736}{3749} and \csource{find_points.c}{3763}{3763}; the listing of the
    denominators of a shape and its test ($\mathrm{squaresDens}$, $\mathrm{squares1Dens}$,
    $\mathrm{wordBits}$, $\mathrm{wordsDens}$, $\mathrm{Shape.dens}$, $\mathrm{Shape.keep}$) becomes
    a definition of the proofs, the node saying that the model has the blocks and no listing (I.11);
    the passage after Lemma~\ref{lem:shape_dens} on the run with its links, its contract and its
    plain equation is on the run over the blocks (I.11); Section~\ref{sec:denSearch}: what the
    search does, with Definition~\ref{def:sievingInfo} and the cut points of the parameters (I.11),
    for well-formed parameters (I.10), and the list of the C ranges of the node of the search, which
    loses lines 2793-2795 (I.8), the lines of the choice and the sieve (I.11) and those of the
    report (I.15); Theorem~\ref{thm:findPoints_order} takes $s$ (I.10) and its proof goes through
    the block-listing lemma (I.11); the deviations that leave the final \verb|pn_lim|, the blocks,
    the reciprocal of the first word of a block and the correction to stage~I (I.11) and the one on
    \verb|RATPOINTS_VERBOSE| (I.15) are reworded, and the closing passage, which leaves the blocks
    with the correction and the number of primes examined to stage~I (I.11); and the files section:
    the shapes and the run in \verb|Ratpoints/Model/Denominators.lean|, the run over a shape in
    \verb|Ratpoints/Proof/Denominators.lean| (I.11) and the contracts from one piece on in
    \verb|Ratpoints/Proof/Sift.lean| (I.4).
  \item Chapter~\ref{chap:tables}: the introduction refers to the options of the sieve for
    Definition~\ref{def:sieveParams} (I.12) and speaks about the choice of the search (I.16); the
    tables of a prime power (Definition~\ref{def:examinePower}) take whether denominators divisible
    by $p$ occur and give the density, with \lname{Ratpoints.Sieve.Power.useful}, their links
    becoming \csource{find_points.c}{2267}{2282}, \csource{find_points.c}{2285}{2293},
    \csource{find_points.c}{2301}{2301}, \csource{find_points.c}{2330}{2330} and
    \csource{find_points.c}{2331}{2331}, their uses gaining Definition~\ref{def:densities}, and
    Lemmas~\ref{lem:power_admits} and~\ref{lem:examinePrime_eq} the forms with the new results, the
    second with \lname{Array.any\_push\_id\_eq\_false\_iff}, the fact on arrays that its contract
    uses (I.2); the sieve entry (Definition~\ref{def:SieveEntry}) has its density and the count of
    its factors (I.2); the loop of the tables of stage~F1 (Definition~\ref{def:examinePrimes} and
    Lemma~\ref{lem:examinePrimes}), the factors of a modulus (Definition~\ref{def:modulusFactors}
    and Lemma~\ref{lem:modulusFactors}), the explicit sieve entry of a modulus
    (Definition~\ref{def:makeModulus}, \verb|makeModulus?|) and the explicit $\mathrm{sievingInfo}$
    (Definition~\ref{def:sievingInfo} and Lemma~\ref{lem:sievingInfo}, renamed with their contract
    at I.8) give their links to the components of the choice (I.3, I.7, I.8) and are deleted (I.12),
    \lname{Ratpoints.getElem!\_primeTable\_le\_of\_le} going to Definition~\ref{def:primeTable},
    \lname{Ratpoints.IsModulusEntry.isSound} to Lemma~\ref{lem:sieveEntry_admits}, and their uses
    and references going to Definitions~\ref{def:tablesLoop}, \ref{def:compositeModuli}
    and~\ref{def:choice} and Lemmas~\ref{lem:compositeModuli} and~\ref{lem:choice}; the rows of a
    modulus (Lemma~\ref{lem:sieveEntry_admits}) keep the factors' soundness and the entries' facts,
    without the invariants of \verb|makeModulus?|, their statement and proof rewritten without
    Lemma~\ref{lem:modulusFactors}, and their uses those that the Lean then implies,
    Definitions~\ref{def:SieveEntry}, \ref{def:examinePrime}, \ref{def:squaresMod}, \ref{def:mulMod}
    and~\ref{def:packingConstants} and Lemma~\ref{lem:mulMod}, the proof
    Lemma~\ref{lem:classOffsets} (I.12); Section~\ref{sec:tablesSearch}, on $\mathrm{Params.WF}$,
    the divisor form and the parameters of the command line, and the paragraph after
    Lemma~\ref{lem:sievingInfo}, on the search for every well-formed choice, speak about the choice
    of the search (I.12); the program and its differential tests (Section~\ref{sec:tablesCli}): the
    program chooses composite moduli itself (I.11); the deviations: the count of the residues (I.2),
    the density, the ranking and the choice not modelled (I.11), those of the explicit route, with
    their references to Definition~\ref{def:makeModulus} (I.12), and the closing passage that leaves
    the density, \verb|add_moduli|, the ranking and the choice to stage~I (I.11), its reference to
    Definition~\ref{def:sievingInfo} going to Definition~\ref{def:choice} (I.12); and the files
    section: the choice in \verb|Ratpoints/Model/Choice.lean| and its contracts in
    \verb|Ratpoints/Proof/Choice.lean| and \verb|Ratpoints/Proof/ChoiceModuli.lean| (I.7, I.8), the
    factors, the explicit entry, the loop of the tables and the well-formedness gone (I.12).
  \item Chapter~\ref{chap:stages}: the introduction's sentence on \verb|wf_sieveParams|, which the
    challenge imports, says that stage~I drops it (I.10); the reductions by a multiplication
    (Definition~\ref{def:mulMod}) refer to Definition~\ref{def:choice} for the array \verb|magics|
    (I.12); the test of a candidate and the extraction (Definitions~\ref{def:accepted}
    and~\ref{def:siftBitArray}, Lemmas~\ref{lem:accepted} and~\ref{lem:siftBitArray}) count, their
    uses gaining Definition~\ref{def:counters}, and the test gains
    \lname{Ratpoints.Sieve.fst\_run\_accepted}, that its verdict does not depend on the counters
    (I.4); the passage after the staged sieve, on the run with the cut points of the parameters and
    what \verb|sieving_info| gives (Lemma~\ref{lem:sievingInfo}), is on the decisions of the blocks
    and the choice (I.11); the section on the cut points (Section~\ref{sec:cuts}) describes the cut
    points of the choice and their correction, not the counting rules of the parameters of the
    executable (I.11); the program and its differential tests (Section~\ref{sec:stagesCli}): the
    program runs the stages of the choice (I.11), and the guards on the stage sizes of the
    parameters go with them (I.12); the deviations that leave the reciprocals of the forbidden
    divisors (I.2), the counts (I.4), the reserve of the reciprocals, the cut points of each block,
    the counting rules and \verb|-R| (I.11) to stage~I are reworded, and the one on the pool of
    threads says that the model has none (I.16); the closing passage that leaves the correction, the
    counters, the choice of the cut points, the pool of threads and the first word of a block to
    stage~I (I.11); and the files section: the run, the parameters of the command line and the
    reciprocals of \verb|sieving_info| (I.11, I.12), the contracts of one piece and one denominator
    (I.4), \verb|wf_sieveParams| (I.12).
  \item Chapter~\ref{chap:packing}: the introduction, on $\mathrm{Params.WF}$ and what the packing
    needs of the parameters (I.12); the passage on the sieve in the packing, the store fitting the
    entries of the first two stages and what \verb|sieving_info| gives
    (Lemma~\ref{lem:sievingInfo}), is on the choice (I.11); what the run needs of the search
    (Lemma~\ref{lem:soundRun}) no longer reads $s_2$, with the entries that count no factor prime
    entries and the store fitting all the entries (I.11); the row offsets of the classes
    (Lemma~\ref{lem:classOffsets}) are made for all the entries, the reserve included (I.11); the
    bits of an interval (Definition~\ref{def:siftBits}, Lemma~\ref{lem:siftBits}) count, their uses
    gaining Definition~\ref{def:counters} (I.4); the deviations: the reserve of the offsets (I.11)
    and the reference to Definition~\ref{def:makeModulus}, which goes to
    Definition~\ref{def:compositeModuli} (I.12); the closing passage that leaves the cost model of
    the packing (I.3) and the reserve of the offsets (I.11) to stage~I; and the files section: the
    run (I.11), the contracts from one piece on (I.4).
  \item Chapter~\ref{chap:rows}: the introduction, on $\mathrm{Params.WF}$ (I.12); the tables of the
    store (Definition~\ref{def:storeTables}) are made as the choice makes them, at its end for all
    the entries, those of the third stage among them (I.11), and refer to
    Definitions~\ref{def:tablesLoop}, \ref{def:choice} and~\ref{def:compositeModuli} for the deleted
    nodes (I.12); the store fits the sieve entries (Lemma~\ref{lem:storeTables}) with the forms with
    the flag of the prime powers, its proof using Lemma~\ref{lem:densities} (I.2), its use of
    Definition~\ref{def:makeModulus} dropped and its references going to
    Definition~\ref{def:compositeModuli} and Lemma~\ref{lem:choice}, its statement about the store
    of the choice (I.12); the passage after Lemma~\ref{lem:buildRows} on the run, the store of
    \verb|sieving_info| and the fit of the first $s_2$ entries is on the blocks and the store of the
    choice, which fits all the entries (I.11); the deviations that leave the acquire loads and the
    threads to stage~I say that the model has no threads (I.16), and their reference to
    Definition~\ref{def:makeModulus} goes (I.12); the closing passage that leaves the lock of the
    tables, the pool of threads and the orderings to stage~I (I.11); and the files section: the
    entries and \verb|sieving_info| in \verb|Ratpoints/Model/FindPoints.lean| (I.8, I.12), the run
    and its contract (I.11), the contracts of one piece and one denominator (I.4).
  \item Chapter~\ref{chap:output}: the introduction: message~4 is printed (I.13), and the paragraph
    on the listing of the denominators, which made a stop pay for the listing until the blocks of
    stage~I, goes (I.11); what the C code does: the numbers of the moduli used, left to stage~I
    (I.13); the delivery (Definition~\ref{def:delivery}) gains \lname{Ratpoints.Delivery.print}
    (I.10); the delivery of the check and the stops of the sieve (Lemmas~\ref{lem:checkDelivery}
    and~\ref{lem:sieveStops}) take the counters and the general list lemmas, the uses of the first
    gaining Definition~\ref{def:counters} and its proof Lemma~\ref{lem:accepted} (I.4); the sieving
    (Definition~\ref{def:sieving}) takes the options and the result of the choice and returns the
    ends $s_2$ and $s_3$ of the last block (I.11); the outcome (Definition~\ref{def:Outcome}) has
    the numbers of the moduli used, its link \csource{ratpoints.h}{377}{383} in place of
    \csource{ratpoints.h}{377}{377} (I.13); the search with a callback
    (Definition~\ref{def:findPointsWith}) takes $s$ and the print channel, its links going to
    Definition~\ref{def:findPointsVerbose}, its uses gaining that node and losing
    Definition~\ref{def:Params} (I.10); the list of the search (Lemmas~\ref{lem:searchList}
    and~\ref{lem:findPoints_searchList}) on the choice, the uses of Lemma~\ref{lem:searchList} going
    to Lemma~\ref{lem:choice} for Lemma~\ref{lem:sievingInfo} and those of its proof gaining
    Theorem~\ref{thm:blocks} and Lemma~\ref{lem:noCheckList}, for the denominators of the search
    without duplicates (I.11), for every $s$ with no hypothesis (I.10); Theorems~\ref{thm:protocol},
    \ref{thm:noCheck} and~\ref{thm:order} take $s$ (I.10), and the order goes through the
    block-listing lemma (I.11); Section~\ref{sec:certifiedG} records the change (I.10); the report
    (Definition~\ref{def:cliReport}) prints message~4, its link \csource{main.c}{979}{1033} in place
    of \csource{main.c}{979}{991} and \csource{main.c}{999}{1033} (I.13), and its uses go to
    Definition~\ref{def:sieveOptions} for Definition~\ref{def:sieveParams} (I.12); the paragraph
    after it, on $\mathrm{sieveParams}(s)$, and the program theorem (Theorem~\ref{thm:cliRun}) take
    the options of the sieve, its uses losing Definition~\ref{def:sieveParams} (I.10), and the
    hypothesis on \verb|-v|, its proof using Lemma~\ref{lem:cliVerbose}, and its
    \lname{Ratpoints.Cli.Data.keeps}, the points that the callback keeps, goes to the callback of
    the program (Definition~\ref{def:cliProcess}) (I.15); the section on the differential tests:
    message~4 compared (I.13), the report under \verb|-v| (I.15), the pairs under \verb|-x| (I.11),
    and the runs skipped for message~4 (I.13) and for the report (I.15); the deviations that leave
    the stop of a thread of the pool, the sieving thread and the blocks to stage~I (I.11), the
    numbers of the moduli used and message~4 (I.13), the report of the library (I.15) and the tuning
    options (I.14) are reworded; the record of the lines left free for stage~I: the tuning fields,
    the defaults of \verb|main.c| and the private data (I.8), the number of threads and line 5163
    (I.11), and the numbers of the moduli used with their lines and message~4 (I.13); and the files
    section: the listing and the run in \verb|Ratpoints/Model/Denominators.lean| and the run on the
    delivery in \verb|Ratpoints/Proof/Denominators.lean| (I.11), \verb|Ratpoints/Proof/Sift.lean|
    (I.4).
  \item Chapter~\ref{chap:region}: the introduction and the description of the C code no longer
    leave the estimate of the run to stage~I (I.3), and the introduction no longer leaves the tuning
    options unread (I.14); the doubles (Definition~\ref{def:Double}) gain
    \lname{Ratpoints.Double.toFloat} (I.3), \lname{Ratpoints.Double.minusOne} (I.8), and
    \lname{Ratpoints.Double.ofFloat} and \lname{Ratpoints.Double.printFixed} (I.15); the passage on
    the options of the search and the deviation that call \verb|sp3| stage~I's give it to
    $\mathrm{Params}$ (I.8); the loop over the intervals (Lemma~\ref{lem:siftIntervals}) counts, its
    uses gaining Definition~\ref{def:counters} (I.4); the passage on what \verb|sieving_info|
    computes for every height refers to Lemma~\ref{lem:choice} for Lemma~\ref{lem:sievingInfo}
    (I.12); the passages of Section~\ref{sec:regionSearch} for well-formed parameters,
    Lemma~\ref{lem:elimSearchRegion} and Theorem~\ref{thm:sturmSpeed} take $s$ (I.10);
    Section~\ref{sec:certifiedH} records the change (I.10); the Sturm step of the search
    (Definition~\ref{def:sturmRegion}) gives its lines of the report to Definition~\ref{def:report},
    its link becoming \csource{find_points.c}{4857}{4859}, \csource{find_points.c}{4862}{4863} and
    \csource{find_points.c}{4879}{4892} (I.15); Theorem~\ref{thm:cliRegion} gains that \verb|-q|
    cancels \verb|-v|, and its uses lose Theorem~\ref{thm:cliRun}, which they reached only for the
    points that the callback keeps (I.15); the section on the differential tests: message~4 compared
    (I.13), the pairs under \verb|-x| compared exactly and the counts of the weaker check (I.11);
    the record of the lines left free after stage~H, each line leaving it with the item that records
    it: the tuning macros (I.1), the estimate (I.3), the counters and the decision of the sieving
    thread (I.4, I.6), the blocks (I.5), the choice, the tuning fields, the private data and the
    defaults of \verb|main.c| (I.8), the correction (I.9), the threads (I.11), and the numbers of
    the moduli used and message~4 (I.13); and the files section: the run (I.11).
  \item The appendix: the generated table follows each item.
\end{itemize}
Outside the blueprint: the description of the repository states the eleven statements (I.10), the
pairs under \verb|-x| compared exactly (I.11) and the tuning options read (I.14);
\verb|formalization.yaml| its statements and their count (I.10) and drops the entry of
\verb|Ratpoints.wf_sieveParams| (I.12); \verb|comparator/README.md| the statements and the names
taken from the model (I.10) with the record of each acceptance (I.10, I.11, I.13, I.15); and the
challenge its module docstring (I.10).

\section{Files and order of work}\label{sec:moduliFiles}

The model, in the folder \verb|Ratpoints/Model|: the constants in \verb|Constants.lean| (about $650$
lines); the cost model in a new file \verb|Costs.lean| (about $700$ lines); the choice in a new file
\verb|Choice.lean| (about $930$ lines); the blocks, the correction and the run in a new file
\verb|Blocks.lean| (about $1010$ lines); the report of the library in a new file \verb|Report.lean|
(about $340$ lines); the densities and the fields of the sieve entries in \verb|Tables.lean|; the
reciprocal of a forbidden divisor in \verb|Denominators.lean|, which loses the listing and the run
(about $880$ lines); the counters in \verb|Candidate.lean| and \verb|Sieve.lean|; the options of the
sieve and $\mathrm{Params}$ in \verb|Params.lean| (about $270$ lines); the search with the choice in
\verb|FindPoints.lean|; the doubles as \verb|Float|s and \verb|%.pf| in \verb|Double.lean|; the
print channel in \verb|Delivery.lean|; and the tuning options, message~4 and the report under
\verb|-v| in \verb|Cli.lean| (about $1230$ lines). The proofs, in the folder \verb|Ratpoints/Proof|,
use Lean core and \verb|Std| only in these files: the pool of the candidates in a new file
\verb|ChoicePool.lean| (about $310$ lines); the counters of the walk, the extraction and the scan in
a new file \verb|ScanCounts.lean| (about $670$ lines); the arithmetic of the counters and the test
of a candidate in \verb|Stages.lean|; \verb|Sift.lean| split into \verb|SiftBits.lean| (the masks,
one piece and the bits of an interval, about $1020$ lines) and \verb|Sift.lean| (one denominator,
about $1090$ lines); the plan of the blocks in a new file \verb|BlockPlan.lean| (about $460$ lines);
the frame of the store in \verb|RowStore.lean| (about $1450$ lines); and \verb|Denominators.lean|,
which keeps only what \verb|run_setup| returns (about $50$ lines). They use Mathlib in these: the
composite moduli and the stages in a new file \verb|ChoiceModuli.lean| (about $780$ lines); the
contract of the choice in a new file \verb|Choice.lean| (about $550$ lines); the loops of the shapes
in a new file \verb|BlockLoops.lean| (about $1100$ lines); one block and the block-listing lemma in
a new file \verb|Blocks.lean| (about $430$ lines); the correction, the run and the block and lag
lemmas in a new file \verb|Correction.lean| (about $1020$ lines); the listing of the denominators as
a definition of the proofs in \verb|Shapes.lean|; the search with the choice in
\verb|FindPoints.lean|, the theorems in \verb|Output.lean| and the program in \verb|Cli.lean|; and
\verb|Moduli.lean|, \verb|Forbidden.lean| and \verb|Params.lean|, which lose the explicit route (to
about $910$, $960$ and $240$ lines). Every Lean file stays below $1500$ lines. The new modules join
their part roots and \verb|Ratpoints.lean|.

One commit per item, each passing the project's checks, the build with other values of the constants
and the differential tests, with the nodes and the passages of the earlier chapters that it changes,
the C paragraphs of the model and the \verb|\csource| links of its nodes, locked, the C links moved
from their old owners and the exclusions in the same commit, the \verb|\lname| of its declarations
turned into \verb|\lean|, and the dependency audit of the blueprint after each item that changes
nodes. The record of the correspondence admits no overlapping ranges, so that a C link moves in the
item that brings its new owner: the lines of the explicit route with I.3, I.7 and I.8 (the paragraph
after Definition~\ref{def:choice}), the lines of the run with I.5, I.6 and I.9, and those of the
search with I.8, I.11 and I.15. The output of \verb|ratpoints-lean| does not change up to I.10; it
changes at I.11, where the pairs under \verb|-x| become those of ratpoints on the runs without
tuning options, at I.13 with message~4, at I.14 with the tuning options and at I.15 with the report
under \verb|-v|; the harness compares each part of the output from the item that makes it that of
ratpoints, on the runs for which it does. The certified statements change once, at I.10, and keep
their text from then on, I.11 changing their proofs; comparator runs with \verb|--paranoid| at I.10,
and plainly at I.11 and at I.13 and I.15, which change definitions that the challenge imports (the
outcome, and the search with a callback through the report).
\begin{enumerate}
  \item[I.1] The constants (Definition~\ref{def:choiceConstants}), \verb|Float|-valued and the
    option, with $\mathrm{blockLen\_pos}$ and the guards of their bits, and $\mathrm{sp2Extra}$,
    $\mathrm{defaultNumPrimes}$ and $\mathrm{defaultMaxForbidden}$ with their links from the
    parameters of the executable; the build with other values of the constants learns \verb|Float|
    constants and the option, with the variants of the lag, the block length and the blocks wanted;
    and the nodes and passages of the earlier chapters that Section~\ref{sec:moduliEarlier} assigns
    to it. Nothing uses the new constants yet.
  \item[I.2] The fields of the entries (Definition~\ref{def:densities}): the code and the density of
    the examination of a prime, the flag and the density of a prime power with its parts and
    $\mathrm{Power.useful}$, the density and the count of a sieve entry, the reciprocal of a
    forbidden divisor, with Lemma~\ref{lem:densities}; the explicit route fills the new fields,
    which nothing reads yet; and the nodes and passages of the earlier chapters that
    Section~\ref{sec:moduliEarlier} assigns to it.
  \item[I.3] The cost model in \verb|Ratpoints/Model/Costs.lean| (Definitions~\ref{def:costModel}
    and~\ref{def:runShape}), with no proofs, and $\mathrm{Double.toFloat}$, which joins
    Definition~\ref{def:Double}; the lines 1619-1623 of the offset of the second stage from the
    parameters of the executable; \verb|ceil|, \verb|floor| and the count of the set bits of a word
    recorded in the \verb|Deviation:| paragraphs of the components that use them; and the nodes and
    passages of the earlier chapters that Section~\ref{sec:moduliEarlier} assigns to it. Nothing
    calls the cost model yet.
  \item[I.4] The counters through the sieve (Definition~\ref{def:counters}) with their closed forms
    (Lemma~\ref{lem:counters}), \verb|Ratpoints/Proof/Sift.lean| split and
    \verb|Ratpoints/Proof/ScanCounts.lean| new; and the nodes and passages of the earlier chapters
    that Section~\ref{sec:moduliEarlier} assigns to it. The run of stage~H discards the counters.
  \item[I.5] The plan of the blocks and the four loops (Definitions~\ref{def:blockPlan}
    and~\ref{def:blockLoops}), with Lemmas~\ref{lem:blockPlan} and~\ref{lem:blockLoops}, the links
    taken from the listing and the run; and the nodes and passages of the earlier chapters that
    Section~\ref{sec:moduliEarlier} assigns to it. Nothing calls them yet.
  \item[I.6] One block (Definitions~\ref{def:runBlock} and~\ref{def:decisionWF},
    Lemma~\ref{lem:runBlock}, with the search for the decision in force), the block-listing lemma
    (Lemma~\ref{lem:blockDens}), stated over the listing of the model, and the frame of the store,
    all of \verb|Ratpoints/Proof/Blocks.lean|, with no unfolding lemma of one block; the links
    3889-3898, 3911, 3913 and 4077-4078 of one block from the run; and the nodes and passages of the
    earlier chapters that Section~\ref{sec:moduliEarlier} assigns to it. Nothing calls them yet.
  \item[I.7] The pool, the composite moduli, the first two stages and the third stage
    (Definitions~\ref{def:compositeModuli}, \ref{def:firstStages} and~\ref{def:thirdStage}, but for
    $\mathrm{thirdOf}$; Lemmas~\ref{lem:choicePool} and~\ref{lem:compositeModuli}), with the C links
    that the paragraph after Definition~\ref{def:choice} gives to I.7; and the nodes and passages of
    the earlier chapters that Section~\ref{sec:moduliEarlier} assigns to it. Nothing calls them yet.
  \item[I.8] The choice (Definitions~\ref{def:tablesLoop} and~\ref{def:choice},
    Lemma~\ref{lem:choice}), $\mathrm{thirdOf}$ with $\mathrm{thirdOf\_spec}$, the new fields of
    $\mathrm{Params}$, and the options of the sieve with their normalization, their tuning and their
    defaults (Definition~\ref{def:sieveOptions}), with $\mathrm{arraySizeAt}$ and
    $\mathrm{Double.minusOne}$; the explicit $\mathrm{sievingInfo}$, \verb|IsSievingInfo| and
    \verb|sievingInfo_spec| renamed to \verb|sievingInfoOf|, \verb|IsSievingInfoOf| and
    \verb|sievingInfoOf_spec|, and the parameters of the executable filling the new fields of
    $\mathrm{Params}$, their docstring naming the fields that the explicit run does not read; the C
    links of the choice from the explicit route and from the search (lines 2794-2795), the private
    fields of the arguments to their owners and the exclusions of those that hold memory; the
    \verb|Deviation:| of the result of the choice corrected: \verb|stage3_list| is read, as the
    buffer of the table of the third stage, which the model replaces by the thunk of the checks; and
    the nodes and passages of the earlier chapters that Section~\ref{sec:moduliEarlier} assigns to
    it. Nothing calls the choice yet.
  \item[I.9] The correction (Definition~\ref{def:correction}), with Lemmas~\ref{lem:correctionWF}
    and~\ref{lem:correctionAfter} and the lag lemma (Theorem~\ref{thm:lag}), which do not need the
    run; the lines 3755-3762 from the shapes and line 3861 from the run; the exclusion of
    \verb|RP_PRIME_STATS|; and the nodes and passages of the earlier chapters that
    Section~\ref{sec:moduliEarlier} assigns to it. Nothing calls the correction yet.
  \item[I.10] The interface with the behaviour unchanged: $\mathrm{findPoints}$ and
    $\mathrm{findPointsWith}$ on $s$ through the parameters of the executable and the explicit
    route, the print channel with $\mathrm{findPointsWork}$ and $\mathrm{findPointsVerbose}$
    (Definition~\ref{def:findPointsVerbose}), $\mathrm{Delivery.print}$ and
    Lemma~\ref{lem:choiceSearch}; the links of the search with a callback and of the program; the
    eleven statements of Section~\ref{sec:certifiedI}, the solution and comparator with
    \verb|--paranoid|, its acceptance recorded; the statements in the description of the repository,
    \verb|formalization.yaml|, \verb|comparator/README.md| and the docstring of the challenge; and
    the nodes and passages of the earlier chapters that Section~\ref{sec:moduliEarlier} assigns to
    it.
  \item[I.11] The switch of the choice and of the run: $\mathrm{chooseAndSearch}$ with the choice
    (Definition~\ref{def:chooseAndSearch}), the run over the blocks with the correction
    (Definition~\ref{def:sieveRun}), the listing moved to the proofs, the block lemma
    (Theorem~\ref{thm:blocks}) and the form with the blocks sifted apart
    (Theorem~\ref{thm:threads}), the order through the block-listing lemma; the route of the proofs
    through $\mathrm{IsSievingInfo}$ and the well-formed decisions, the statements unchanged, and
    the docstrings of the contracts of the search and of its run, which say that the test does not
    depend on the search intervals, corrected; comparator, its acceptance recorded; the exclusions
    of the threads; the sentence of the description of the repository on \verb|-x|; the harness: the
    pairs under \verb|-x| compared exactly on the runs without \verb|-r -R -U -C -A -Q -W|, the
    suite \verb|adapt| on its $20$ command lines without tuning options, without \verb|-A| and with
    \verb|-A 1|, the new runs of the high suite, those with \verb|-A 2| compared exactly from I.14,
    and the pair with the block length $3$ on the runs without tuning options; and the nodes and
    passages of the earlier chapters that Section~\ref{sec:moduliEarlier} assigns to it.
  \item[I.12] The deletion of the explicit route and its proofs: the parameters of the executable,
    the loop of the tables of stage~F1, the factors of a modulus, \verb|makeModulus?|, the explicit
    $\mathrm{sievingInfo}$ with its contract, and $\mathrm{Params.WF}$ with its instance, its lemmas
    and \verb|getElem!_primeTable_inj|; \verb|getElem!_primeTable_le_of_le| and
    \verb|IsModulusEntry.isSound| moved; the docstrings of the model that name $\mathrm{Params.WF}$,
    in \verb|Ratpoints/Model/Candidate.lean| and \verb|Ratpoints/Model/Params.lean|, reworded; the
    nodes that describe these, and every reference and use of a deleted label in
    Chapters~\ref{chap:spec} to~\ref{chap:moduli} retargeted or dropped, as
    Section~\ref{sec:moduliEarlier} lists them, the references of this chapter turned into plain
    names; the entry of \verb|wf_sieveParams| in \verb|formalization.yaml|; and the nodes and
    passages of the earlier chapters that Section~\ref{sec:moduliEarlier} assigns to it.
  \item[I.13] The numbers of the moduli used in the outcome and message~4 in the program, with their
    links; comparator, its acceptance recorded; the harness compares message~4 on the runs without
    tuning options; and the nodes and passages of the earlier chapters that
    Section~\ref{sec:moduliEarlier} assigns to it.
  \item[I.14] The tuning options of the program (Definition~\ref{def:cliTuning}, the reading of the
    options and its lemmas extended to their branches); the pairs under \verb|-x| and message~4
    compared on every run, the whole suite \verb|adapt|, the runs of the high suite with \verb|-A 2|
    compared exactly, the suites \verb|tuning| and \verb|threads|, and the pair with the block
    length $3$ on every run; the sentence of the description of the repository on the options that
    the program ignored; and the nodes and passages of the earlier chapters that
    Section~\ref{sec:moduliEarlier} assigns to it.
  \item[I.15] The report of the library (Definition~\ref{def:report}) in
    \verb|Ratpoints/Model/Report.lean|, with \verb|%.pf|, $\mathrm{Double.ofFloat}$, the print
    channel of the program and the links of the report from the components that run around it, and
    the program theorem with the hypothesis on \verb|-v| (Lemma~\ref{lem:cliVerbose}); comparator,
    its acceptance recorded; the harness compares the report under \verb|-v| byte for byte and has
    the suite \verb|verbose|; and the nodes and passages of the earlier chapters that
    Section~\ref{sec:moduliEarlier} assigns to it.
  \item[I.16] The texts: the description of the repository, \verb|formalization.yaml| and
    \verb|comparator/README.md| beyond the items, the introduction of the blueprint, and the
    passages of the earlier chapters that stage~I makes untrue beyond those of the items, which
    Section~\ref{sec:moduliEarlier} assigns to it; the check of the correspondence switched to
    \verb|--complete| in CI and in the checks of each commit, and \verb|correspondence/README.md|
    reworded; the generated table.
\end{enumerate}
The blocks come before the choice, since the contract of the choice states that its first decision
is well formed (Definition~\ref{def:decisionWF}, I.6); the correction comes after the choice, since
it reads the fields of the options and of the parameters that the choice brings; and the cost model
comes before both. The switch of the run and the switch of the choice are one item: a run over the
blocks with the explicit choice would need the contract of the choice for the explicit route, a
proof that I.12 would delete again.

Names that the prototype has and the audit of the stage is to review: the two definitions named
$\mathrm{ForbFits}$, \lname{Ratpoints.Sieve.ForbFits} (on lists) and
\lname{Ratpoints.Sieve.Shape.ForbFits}; the names in the style of abbreviations
\lname{Ratpoints.AllQ}, \lname{Ratpoints.TailNf}, \lname{Ratpoints.CandOK},
\lname{Ratpoints.PrimeSeOK}, \lname{Ratpoints.CodeOK}, \lname{Ratpoints.EntryOK},
\lname{Ratpoints.EntriesOK} and the structures named \verb|OK|; the general fact
\lname{Ratpoints.Sieve.Nat.ite\_mod\_add\_one\_eq\_add\_one\_mod}, which landed in the namespace
\verb|Ratpoints.Sieve|, and \lname{Ratpoints.Sieve.take\_append\_take\_drop}; the names
\lname{Ratpoints.Sieve.pre\_siftBits\_counters}, \lname{Ratpoints.Sieve.ite\_lt\_rbaLength\_and} and
\lname{Ratpoints.Cli.input\_eq\_ok}, which do not say what they state; the general facts on arrays,
options and lists of the new proof files that may belong in \verb|Ratpoints/Proof/Containers.lean|
(\lname{Array.Extends}, \lname{Option.eq\_none\_of\_forall\_prod4},
\lname{List.forall\_mem\_drop\_of\_forall\_getElem!} and their kin); the ten to twelve fixed
arguments of \lname{Ratpoints.Sieve.blockReport} and \lname{Ratpoints.Sieve.correctionAfter}, which
a structure could bundle; and whether \verb|Ratpoints/Model/Blocks.lean| splits off the correction
as a file of its own. The names of this chapter are those of the prototype; the audit may rename
them.

\section{The evidence}\label{sec:moduli-evidence}

\begin{itemize}
  \item A reading of the C code of the choice, of \verb|run_shape| and its helpers, of the blocks,
    the counters, \verb|run_fold| and \verb|adapt_primes|, of the threads, of the floating-point
    environment of the build, of the numbers of the moduli used, message~4 and the report under
    \verb|-v|, at the pin, by three readers, each with a second reader who checked the findings: a
    model of the choice in plain IEEE doubles in the order of the C code, with glibc's \verb|log|
    and \verb|pow|, reproduced, from the estimates, the densities and the candidates that a traced
    build of the C program printed, the costs, the keys, the numbers of the moduli of the stages and
    the primes of the third stage on $5901$ runs, the suites, the moduli pinned, the curves of
    \verb|rptest| at the heights $200$, $16383$ and $50000$ and sweeps of the tuning options, and,
    from the traced counters, the decisions of all $2636$ calls of \verb|adapt_primes| (the sketch
    below reproduced the whole choice); the C program with \verb|-t 4|, \verb|-t 0| and \verb|-t 2|
    printed the output of the sequential run on $2351$ sieved runs of the suites and $450$ runs of
    stress, also in its traced decisions and the numbers of the moduli used; glibc's \verb|qsort|
    was stable in $50000$ random trials with ties; \verb|log| and \verb|pow| were measured against
    the correctly rounded values; the build compiles no fused multiply-add. The reading found no
    defect of the behaviour of the C code, and three cosmetic points, recorded with the C code for a
    release after 3.1.1: the sort by density from the end of the second stage,
    \verb|RATPOINTS_DEFAULT_SP1| read nowhere, and two comments that say that \verb|fill_bp_list|
    makes the correction, which \verb|run_fold| makes. A reading of the Lean side checked that
    Lean's \verb|Float| computes the doubles of the cost model as the C program does on $6000$
    inputs, $16$ doubles of the cost model each, compiled and interpreted, and proved, before the
    stage, that a choice driven by \verb|Float| keys is well formed without any fact about
    \verb|Float|, that the cut points of each block keep a sound test, and the block lemma with its
    threaded form.
  \item Before the stage, a sketch of the model in Lean in three parts, the choice, the blocks and
    the report, checked against an instrumented build of the C program that prints every double of
    the choice in hexadecimal: $2349$ command lines of the suites that reach the choice gave
    $144919$ identical lines of the trace, with $2463$ steps of the loop of the tables that look at
    one more prime, $6$ exits in that loop and $1$ in the third stage, composite moduli in $48$ runs
    and a reserve in $2077$, and the $398$ runs of a file of vectors; $8000$ runs with random tuning
    options over three seeds, byte for byte; every fold of the blocks of the runs of \verb|test4.sh|
    at $30000$ to $100000$ ($13201$ lines of the trace, $15$ decisions, $5$ of them taking effect)
    and of $2356$ runs of \verb|rptest| at $16383$ and $50000$ ($188558$ lines, $149$ decisions, in
    $37$ runs taking effect); the reports under \verb|-v| of the $33$ runs \verb|v()| of
    \verb|test4.sh|, $14$ larger runs and $2000$ random runs; and \verb|%.pf| against glibc on
    $3546$ doubles; none failing.
  \item Reviews of the sketch, of the choice against the C code, of the blocks, the counters and the
    report against the C code, of the harness and of the specification, each finding checked by a
    second reader: $50$ findings, $38$ confirmed and $12$ in part, none refuted, no defect of the
    behaviour of the model. They made the C links exact without overlaps, the explicit
    implementations that the sketch kept beside the faithful ones go from the model, the guards
    independent of the platform's \verb|log| and \verb|pow|, the dispatch of the shapes a function
    of its own, the correction of the third stage compute its costs itself as the C code does, and
    gave the harness the runs at $16383$ where the correction takes effect, the comparison of the
    pairs of \verb|test4.sh| with the live C program and the refusal of a \verb|tuning.mk|; the
    sketch with these fixes compared $5397$ times in the default suites, $13$ in \verb|test4high|
    and $92$ in \verb|adapt|, and the pair with the block length $3$ compared $3601$ times in twelve
    default suites and $92$ times in \verb|adapt|, none failing.
  \item A prototype of all the proofs on the sketch: the whole library, the program, the challenge
    and the solution build, with no \verb|sorry| outside the challenge, every proof at most $20$
    lines with no exception of the style lint, standard axioms; comparator with \verb|--paranoid|
    accepts the eleven statements with all six kernels (2026-10-25); the build with other values of
    the constants passes its eleven variants, among them the lag $1$ with the block length $3$ and
    one block wanted; the environment linters pass; the differential tests compare as in
    Section~\ref{sec:moduliTests}; every Lean file is below $1500$ lines; it adds $566$ public names
    and removes $92$. It changed the model of the sketch not at all, and it showed which hypotheses
    the contracts need: the listing of the search without duplicates and the forbidden divisors with
    their reciprocals for the contract of the run per denominator, normalized options for the
    choice, and $\mathrm{IsSievingInfo}$ and the well-formed decisions in place of
    $\mathrm{Params.WF}$.
\end{itemize}

